\documentclass[11pt,a4paper]{article}

\usepackage[utf8]{inputenc}
\usepackage[T1]{fontenc}
\usepackage[italian]{babel}
\usepackage{amsmath,amssymb}
\usepackage{booktabs}
\usepackage[margin=2.5cm]{geometry}

% Sezioni e sottosezioni sono numerate; i titoli di livello inferiore no
\setcounter{secnumdepth}{2}

\newcommand{\file}[1]{\texttt{\detokenize{#1}}}
\newcommand{\Tc}{T_c}
\newcommand{\nom}{\mathrm{nom}}
\newcommand{\lib}{\mathrm{lib}}
\newcommand{\stab}{s^{\mathrm{tab}}}        % punti di SOC delle tabelle

\title{Linearizzazione della predizione\\nell'MPC di bilanciamento attivo}
\author{Modello Simulink \file{Bilanciamento_MPC.slx}}
\date{}

\begin{document}
\maketitle

\noindent
Questo documento descrive come il sottosistema \file{Controllo} del modello
Simulink trasforma il problema di controllo ottimo del paper di Chen, Behal e
Li (IEEE CDC 2021) in un problema quadratico, e in particolare come calcola le
derivate che servono per farlo.

\section{Perché serve linearizzare}

A ogni passo di controllo l'MPC risolve il problema (eq.~6): scegliere le
correnti di bilanciamento $u^1_k,\dots,u^N_k$, costanti su tutto l'orizzonte,
che minimizzano un costo costruito sulle tensioni predette delle celle. Le
tensioni predette dipendono da $u_k$ attraverso il modello di cella, che non è
lineare: i parametri $V_{oc}$, $R_o$, $R_p$, $C_p$ dipendono dallo stato di
carica.

Se il modello fosse lineare, le tensioni predette sarebbero funzioni lineari
di $u_k$: il costo sarebbe quadratico, i vincoli lineari, e il problema un
programma quadratico (QP), per cui esistono solutori rapidi e affidabili come
\texttt{quadprog}. Con il modello non lineare questo non vale più. Il paper
(Remark~3) osserva però che basta linearizzare il modello di predizione a ogni
passo perché il problema torni a essere un QP. È quello che fa il modello
Simulink, in due blocchi: \file{Predizione_linearizzata} calcola la predizione
e le sue derivate, \file{Problema_QP} le inserisce in costo e vincoli.

L'idea è questa. A ogni passo si calcola come evolverebbero le celle se le
correnti di bilanciamento restassero quelle del passo precedente, e di quanto
questa evoluzione cambierebbe per ogni ampere in più. Con queste due
informazioni la tensione predetta di ogni cella è approssimata da una retta
nella sua corrente di bilanciamento, e il problema diventa un QP.

Il documento segue lo stesso percorso:
\begin{itemize}
  \item la sezione~\ref{sec:modello} presenta il modello di predizione, prima
        a tempo continuo e poi nella forma discreta usata dall'MPC;
  \item la sezione~\ref{sec:punto} fissa il punto attorno a cui si linearizza;
  \item la sezione~\ref{sec:sviluppo} scrive la retta e individua i
        coefficienti che restano da calcolare;
  \item la sezione~\ref{sec:derivate} li calcola: è la parte centrale del
        documento;
  \item la sezione~\ref{sec:qp} li inserisce nel costo e nei vincoli del QP;
  \item le ultime tre sezioni riepilogano un passo di controllo, discutono
        l'accuratezza e indicano dove si trova ogni passaggio nel codice.
\end{itemize}

\paragraph{Convenzioni.} In tutto il documento valgono le convenzioni
seguenti.
\begin{itemize}
  \item \emph{Numeri di equazione.} ``eq.~6'' e ``Remark~3'' rimandano al
        paper. I numeri tra sole parentesi, come (3), sono le equazioni di
        questo documento che il codice applica, tutte a tempo discreto. Le
        equazioni a tempo continuo, che servono a capire da dove vengono le
        prime, sono numerate (C1), (C2), \dots; la (T) è la lettura delle
        tabelle; (U1)--(U3) descrivono la corrente di bilanciamento lungo
        l'orizzonte, all'inizio della sezione~\ref{sec:qp}; (Q1)--(Q4) sono i
        passaggi intermedi che portano alla (11), nella
        sezione~\ref{sec:costo}.
  \item \emph{Indici.} $n = 1,\dots,N$ è la cella, scritta come indice in alto
        ($u^n_k$, $I^n$) e sottintesa quando si parla di una cella sola; $k$ è
        l'istante attuale; $j = 1,\dots,p$ è il passo lungo l'orizzonte; $t$ è
        il tempo continuo, misurato a partire dall'istante attuale.
  \item \emph{Correnti di bilanciamento.} $u^n_k$ è la corrente di
        bilanciamento della cella $n$ che l'MPC sceglie all'istante $k$, e
        $u_k = (u^1_k, \dots, u^N_k)$ il vettore delle $N$ correnti; nei
        passaggi su una cella sola, dove $n$ è sottinteso, $u_k$ indica la
        corrente di quella cella. Il pedice è $k$ e non $k+j$ perché la
        corrente scelta in $k$ resta la stessa per tutti i passi
        dell'orizzonte (sezione~\ref{sec:modello}; in forma matematica,
        \eqref{eq:u-costante} nella sezione~\ref{sec:qp}). La corrente scelta
        al passo precedente è $u^n_{k-1}$.
  \item \emph{Derivate.} Il punto indica la derivata rispetto al tempo,
        $\dot V_p = dV_p/dt$; l'apice la derivata rispetto al SOC,
        $\tau' = d\tau/ds$.
  \item \emph{Barra.} Una grandezza con la barra, come $\bar u$ o $\bar Y$, è
        calcolata nel punto di linearizzazione (sezione~\ref{sec:punto}).
  \item \emph{Correnti.} Sono positive in scarica.
\end{itemize}

\section{Modello di predizione}\label{sec:modello}

Sia $k$ l'istante attuale, $\Tc$ il passo di controllo e $p$ l'orizzonte, cioè
il numero di passi su cui si predice. Per la cella $n$ lo stato è la coppia
$\zeta = (s, V_p)$, come nel paper: stato di carica (SOC) e tensione di
polarizzazione. La corrente
della cella è la somma della corrente di pacco e di quella di bilanciamento
(eq.~3),
\begin{equation}
  I^n = i_k + u^n_k ,
\end{equation}
dove $i_k$ è la corrente di pacco misurata all'istante $k$ e $u^n_k$ la
corrente di bilanciamento che l'MPC sceglie all'istante $k$. Sull'orizzonte di
predizione i due addendi sono entrambi costanti: $i_k$ perché della corrente
di pacco non si ha un'anteprima, $u^n_k$ perché così è formulato il problema.
Il paper sceglie infatti una sola corrente per cella e la tiene invariata per
tutti i $p$ passi (eq.~6b, dove compare $i_{k+j} + u^n_k$): per questo il suo
pedice è $k$, l'istante in cui è decisa, e non $k+j$. All'istante successivo
il problema si risolve di nuovo e si sceglie $u^n_{k+1}$.
Lungo l'orizzonte $I^n$ è quindi una costante: non un segnale che varia nel
tempo, ma un \emph{parametro} da cui dipende tutta la traiettoria della cella.
Su questa osservazione si regge il calcolo delle derivate della
sezione~\ref{sec:derivate}.

Il modello è presentato prima a tempo continuo, come equazioni differenziali, e
poi nella forma discretizzata che l'MPC usa davvero.

\subsection{A tempo continuo}\label{sec:continuo}

La cella è descritta da un circuito equivalente: un generatore $V_{oc}$ in
serie a una resistenza $R_o$ e a un ramo formato da $R_p$ e $C_p$ in parallelo.
Detto $t$ il tempo trascorso dall'istante attuale, lo stato evolve secondo due
equazioni differenziali, a cui si aggiunge un'equazione algebrica per l'uscita,
la tensione ai morsetti $y$:
\begin{align}
  \dot s(t) &= -\frac{\eta}{3600\,C^n}\, I^n ,
     \tag{C1}\label{eq:c-soc}\\
  \dot V_p(t) &= -\frac{V_p(t)}{\tau\bigl(s(t)\bigr)}
                 + \frac{I^n}{C_p\bigl(s(t)\bigr)} ,
     \tag{C2}\label{eq:c-vp}\\
  y(t) &= V_{oc}\bigl(s(t)\bigr) - V_p(t) - R_o\bigl(s(t)\bigr)\, I^n ,
     \tag{C3}\label{eq:c-y}
\end{align}
con $\tau(s) = R_p(s)\,C_p(s)$ e condizioni iniziali $s(0) = s_k$,
$V_p(0) = V_{p,k}$, lo stato misurato.
\begin{itemize}
  \item La \eqref{eq:c-soc} è il conteggio della carica: la corrente, positiva
        in scarica, svuota la cella. $C^n$ è la capacità della cella in Ah,
        cioè la carica che può erogare (il 3600 la porta in ampere-secondo),
        ed $\eta$ è il rendimento di carica. $C^n$ non va confusa con $C_p$,
        il condensatore del ramo RC, che si misura in farad. Una cella piena
        non si carica oltre: se $s = 1$ e $I^n < 0$, allora $\dot s = 0$.
  \item La \eqref{eq:c-vp} è il bilancio di corrente del ramo RC,
        $C_p \dot V_p = I^n - V_p / R_p$, diviso per $C_p$. A corrente
        costante $V_p$ tende a $R_p I^n$ con costante di tempo $\tau$.
  \item La \eqref{eq:c-y} è la tensione ai morsetti: quella a vuoto meno le
        cadute sul ramo RC e sulla resistenza serie.
\end{itemize}
In forma compatta, raccogliendo lo stato in $\zeta = (s, V_p)$:
\begin{equation}
  \dot \zeta(t) = f\bigl(\zeta(t),\, I^n\bigr), \qquad
  y(t) = g\bigl(\zeta(t),\, I^n\bigr), \qquad \zeta(0) = \zeta_k ,
  \tag{C4}\label{eq:c-compatta}
\end{equation}
dove $f$ ha per componenti i secondi membri della \eqref{eq:c-soc} e della
\eqref{eq:c-vp}, e $g$ è il secondo membro della \eqref{eq:c-y}. Il modello non
è lineare perché $\tau$, $C_p$, $V_{oc}$ e $R_o$ dipendono dallo stato di
carica, nel modo descritto qui sotto.

\subsection{I parametri in funzione del SOC}\label{sec:tabelle}

Per ciascuno dei quattro parametri $V_{oc}$, $R_o$, $R_p$, $C_p$ è data una
tabella \emph{nominale}, uguale per tutte le celle, con il valore del
parametro in 7 punti di SOC (nel progetto 0; 0,1; 0,25; 0,5; 0,75; 0,9; 1).
Tra un punto e l'altro si interpola linearmente. Le celle vere però non sono
identiche: ognuna ha, per ogni parametro, un \emph{fattore di deviazione} che
moltiplica la tabella nominale. Per esempio la resistenza serie della cella
$n$ è $R_o^n(s) = \varphi^n_{R_o}\,R_o^{\nom}(s)$.

In formule, sia $P$ uno dei quattro parametri, $\stab_1 < \dots < \stab_7$ i
punti di SOC della tabella e $P^{\nom}_1, \dots, P^{\nom}_7$ i valori nominali
corrispondenti. Se $s$ cade nel tratto $\ell$, cioè
$\stab_\ell \le s < \stab_{\ell+1}$, valore e derivata del parametro sono
\begin{equation}
  P(s) = \varphi_P\,\bigl[P^{\nom}_\ell + m_\ell\,(s - \stab_\ell)\bigr],
  \qquad
  P'(s) = \varphi_P\, m_\ell ,
  \qquad
  m_\ell = \frac{P^{\nom}_{\ell+1} - P^{\nom}_\ell}{\stab_{\ell+1} - \stab_\ell} ,
  \tag{T}\label{eq:tabella}
\end{equation}
dove $m_\ell$ è la pendenza del tratto $\ell$ della tabella nominale. Ne seguono tre
fatti, che serviranno nel calcolo delle derivate.
\begin{itemize}
  \item Ogni parametro è una funzione \emph{lineare a tratti} del SOC: su
        ciascun tratto la sua derivata rispetto a $s$ è una costante.
  \item La derivata non richiede alcun calcolo numerico approssimato: è un
        numero che si legge dalla stessa tabella del valore, e che cambia solo
        quando $s$ passa da un tratto all'altro.
  \item Nei casi limite si procede così. Un punto della tabella, dove la
        funzione fa uno spigolo, è assegnato al tratto alla sua destra, come
        dice la condizione $\stab_\ell \le s < \stab_{\ell+1}$; l'ultimo punto
        all'ultimo tratto. Fuori dalla tabella si usa il valore dell'estremo,
        con derivata nulla.
\end{itemize}
Per esempio, per la resistenza $R_p$ di una cella nominale ($\varphi_{R_p} = 1$) in
$s = 0{,}6$: il punto cade nel tratto tra $0{,}5$ e $0{,}75$, ai cui estremi
la tabella vale $0{,}0016~\Omega$ e $0{,}0023~\Omega$. Quindi
\begin{align*}
  m &= \frac{0{,}0023 - 0{,}0016}{0{,}75 - 0{,}5} = 0{,}0028~\Omega , \\[2pt]
  R_p &= 0{,}0016 + 0{,}0028 \cdot (0{,}6 - 0{,}5) = 0{,}00188~\Omega ,
  \qquad
  R_p' = m = 0{,}0028~\Omega .
\end{align*}
La derivata è riferita all'intera escursione del SOC, da 0 a 1:
$R_p' = 0{,}0028~\Omega$ significa che, se il SOC sale di $0{,}01$, $R_p$ sale
di $0{,}000028~\Omega$.

Nel seguito $V_{oc}$, $R_o$, $R_p$, $C_p$ e $\tau = R_p\,C_p$ indicano sempre i
parametri della cella considerata, fattore compreso: ogni cella ha i suoi. Nel
codice valore e derivata sono restituiti insieme da una sola funzione
(\file{qp_tabella.m}; \texttt{leggi\_tabella} in
\file{predizione_linearizzata_semplice.m}).

\subsection{Discretizzazione temporale}\label{sec:discreto}

Il controllo lavora a passi di durata $\Tc$ e gli servono solo i valori dello
stato agli istanti $k+1, \dots, k+p$. Si passa quindi dalle equazioni
differenziali a una regola che fa avanzare lo stato di un passo alla volta. Con
il metodo di Eulero in avanti la derivata è sostituita dal rapporto
incrementale su un passo,
\[
  \dot \zeta \;\simeq\; \frac{\zeta_{k+j} - \zeta_{k+j-1}}{\Tc}
  \qquad\Longrightarrow\qquad
  \zeta_{k+j} = \zeta_{k+j-1} + \Tc\, f\bigl(\zeta_{k+j-1},\, I^n\bigr) ,
\]
con il secondo membro valutato nello stato di \emph{inizio} passo: lo stato
avanza come se la sua derivata restasse, per tutto il passo, quella che ha
all'inizio. Applicando la regola alle \eqref{eq:c-soc}--\eqref{eq:c-vp}, per
$j = 1,\dots,p$ si ottiene (eq.~2a--2c):
\begin{align}
  s_{k+j} &= \min\bigl(s_{k+j-1} - \kappa^n I^n,\; 1\bigr), \label{eq:soc}\\
  V_{p,k+j} &= V_{p,k+j-1} + \Tc\left(-\frac{V_{p,k+j-1}}{\tau(s_{k+j-1})}
              + \frac{I^n}{C_p(s_{k+j-1})}\right), \label{eq:vp}\\
  y_{k+j} &= V_{oc}(s_{k+j}) - V_{p,k+j} - R_o(s_{k+j})\, I^n , \label{eq:y}
\end{align}
con
\begin{equation}
  \tau(s) = R_p(s)\,C_p(s), \qquad
  \kappa^n = \frac{\eta\,\Tc}{3600\,C^n} .
\end{equation}
Tre osservazioni, una per equazione.
\begin{itemize}
  \item Nella \eqref{eq:soc}, $\kappa^n$ è la frazione di carica persa in un
        passo per ogni ampere: è il coefficiente della \eqref{eq:c-soc}
        moltiplicato per $\Tc$. Il minimo impedisce a una cella piena di
        caricarsi oltre.
  \item Nella \eqref{eq:vp} i parametri del ramo RC sono letti nel SOC di
        \emph{inizio} passo, $s_{k+j-1}$, come vuole la regola di Eulero.
  \item La \eqref{eq:y} non contiene derivate rispetto al tempo, quindi non
        c'è nulla da discretizzare: è la \eqref{eq:c-y} calcolata nello stato
        di \emph{fine} passo, perché $y_{k+j}$ è la tensione all'istante
        $k+j$. I parametri $V_{oc}$ e $R_o$ sono quindi letti in $s_{k+j}$.
\end{itemize}
La distinzione tra inizio e fine passo tornerà nel calcolo delle derivate.

\paragraph{Quanto è fedele il modello discreto.} Per il SOC la
discretizzazione è esatta, perché il secondo membro della \eqref{eq:c-soc} è
costante. Per $V_p$ è un'approssimazione. Raccogliendo $V_{p,k+j-1}$ nella
\eqref{eq:vp},
\[
  V_{p,k+j} = \left(1 - \frac{\Tc}{\tau}\right) V_{p,k+j-1}
              + \frac{\Tc}{C_p}\, I^n ,
\]
si vede che a ogni passo la tensione precedente è moltiplicata per
$1 - \Tc/\tau$. Il modello continuo, a parametri costanti, la moltiplicherebbe
invece per $e^{-\Tc/\tau}$: i due fattori sono vicini solo se $\Tc$ è piccolo
rispetto a $\tau$. Inoltre il modello discreto è stabile solo se
$\Tc < 2\tau$, e non oscilla solo se $\Tc < \tau$. Con i valori di
\file{init_parametri.m} ($\Tc = 20$~s, $\tau$ tra circa 30 e 106~s nei punti
delle tabelle nominali) le due condizioni sono rispettate, ma dove $\tau$ è
minima $\Tc/\tau$ vale circa 0,67: lì il fattore discreto è 0,33 e quello
continuo 0,51. Il valore di regime, $R_p I^n$, è lo stesso nei due modelli;
cambia la rapidità del transitorio.

In ogni caso il modello di predizione dell'MPC è quello discreto
\eqref{eq:soc}--\eqref{eq:y}: è su questo che si calcolano costo, vincoli e
derivate.

\section{Punto di linearizzazione}\label{sec:punto}

Linearizzare significa sostituire una curva con la sua retta tangente, e una
tangente si traccia in un punto preciso. Qui il punto è fissato da tre dati,
tutti noti all'istante $k$:
\begin{itemize}
  \item lo stato misurato $(s_k, V_{p,k})$ di ogni cella;
  \item la corrente di pacco misurata $i_k$;
  \item la mossa del passo precedente, $\bar u^n = u^n_{k-1}$.
\end{itemize}
Ponendo $u^n_k = \bar u^n$, cioè con la corrente $I^n = i_k + \bar u^n$, le
\eqref{eq:soc}--\eqref{eq:y} danno, per ogni cella, la \emph{traiettoria di
linearizzazione}
\[
  \bar Y_{j,n} = y^n_{k+j}(\bar u^n), \qquad
  \bar S_{j,n} = s^n_{k+j}(\bar u^n), \qquad j=1,\dots,p,\; n=1,\dots,N .
\]
È la predizione che si avrebbe se le correnti di bilanciamento restassero
quelle del passo precedente. Ogni cella la percorre con i \emph{propri}
parametri, cioè con i suoi fattori di deviazione e la sua capacità: la barra
indica il punto di linearizzazione, non la cella nominale.

La scelta di $u_{k-1}$ è naturale perché le correnti di bilanciamento variano
lentamente: la soluzione del passo attuale è vicina a quella precedente, e lì
l'approssimazione lineare è più accurata.

Stato misurato, corrente di pacco e mossa precedente cambiano a ogni passo di
controllo, e con loro cambia la traiettoria. Non si linearizza quindi una volta
per tutte attorno a un punto di equilibrio: si linearizza attorno a una
traiettoria, e lo si rifà a ogni passo. È lo schema che in letteratura va sotto
il nome di linearizzazione successiva.

\section{Sviluppo al primo ordine}\label{sec:sviluppo}

Le celle non si influenzano tra loro nella predizione: tensione e SOC della
cella $n$ dipendono solo dalla sua corrente $u^n_k$. Per ogni cella e per ogni
passo $j$ basta quindi lo sviluppo al primo ordine di una funzione di una sola
variabile, attorno a $\bar u^n$:
\begin{equation}
\begin{aligned}
  y^n_{k+j}(u^n_k) &\;\simeq\; \bar Y_{j,n} + d_{j,n}\,\bigl(u^n_k - \bar u^n\bigr),
  &\qquad
  d_{j,n} &= \left.\frac{\partial y^n_{k+j}}{\partial u^n_k}\right|_{\bar u^n} , \\[6pt]
  s^n_{k+j}(u^n_k) &\;\simeq\; \bar S_{j,n} + \alpha_{j,n}\,\bigl(u^n_k - \bar u^n\bigr),
  &\qquad
  \alpha_{j,n} &= \left.\frac{\partial s^n_{k+j}}{\partial u^n_k}\right|_{\bar u^n} .
\end{aligned}
\label{eq:lin}
\end{equation}
Gli istanti nei pedici sono due e hanno ruoli diversi: $k+j$ è l'istante in
cui si guarda la cella, $k$ quello in cui si sceglie la corrente. La stessa
$u^n_k$ agisce su tutti i passi dell'orizzonte, quindi la tensione e il SOC a
ogni passo $j$ sono funzioni di questo solo numero.
I termini noti $\bar Y$ e $\bar S$ sono la traiettoria di linearizzazione
della sezione~\ref{sec:punto}. $d_{j,n}$ e $\alpha_{j,n}$ sono i \emph{coefficienti
di sensibilità} della predizione: dicono di quanto cambiano, tra $j$ passi, la
tensione e il SOC della cella $n$ per ogni ampere in più di corrente di
bilanciamento. Si scrivono in minuscolo come la loro versione a tempo
continuo, $d(t)$ e $\alpha(t)$, che compare nella sezione~\ref{sec:derivate}.
Per ogni $j$ e $n$ la prima riga della \eqref{eq:lin} è la
retta tangente, nel punto $\bar u^n$, alla curva vera $y^n_{k+j}(u^n_k)$:
coincide con essa in $\bar u^n$ e se ne allontana man mano che $u^n_k$ si
scosta da $\bar u^n$.

Poiché ogni cella dipende solo dalla propria corrente, bastano quattro matrici
$p\times N$ (nel codice \file{Y_bar}, \file{S_bar}, \texttt{d}, \texttt{alfa})
invece di una jacobiana $pN \times N$.

Se il modello fosse lineare, $d$ e $\alpha$ sarebbero numeri fissi, calcolabili una
volta per tutte, e la \eqref{eq:lin} sarebbe un'uguaglianza. Con il modello non
lineare dipendono invece dal punto di linearizzazione: vanno ricalcolati a ogni
passo di controllo, e la \eqref{eq:lin} è un'approssimazione, buona finché
$u_k$ resta vicina a $\bar u$.

A questo punto $\bar Y$ e $\bar S$ si ottengono facendo avanzare il modello
\eqref{eq:soc}--\eqref{eq:y} con la corrente $i_k + \bar u^n$. Restano da
calcolare $d$ e $\alpha$: è l'argomento della prossima sezione.

\section{Calcolo delle derivate}\label{sec:derivate}

Questa sezione calcola i coefficienti di sensibilità $d_{j,n}$ e $\alpha_{j,n}$
della \eqref{eq:lin}. Conviene ricordare che cosa sono: le derivate parziali
della tensione e del SOC predetti, all'istante $k+j$, rispetto alla corrente
di bilanciamento della stessa cella, valutate nel punto di linearizzazione.
Tre premesse valgono per tutta la sezione.
\begin{itemize}
  \item \emph{Si deriva rispetto a $I$.} Poiché $I^n = i_k + u^n_k$ e $i_k$ è
        fissata, un ampere in più di $u^n_k$ è un ampere in più di $I^n$:
        derivare rispetto a $u^n_k$ equivale a derivare rispetto a $I^n$. Il
        punto di linearizzazione $u^n_k = \bar u^n$ corrisponde alla corrente
        $I^n = i_k + \bar u^n$. Le formule che seguono valgono per qualunque
        valore della corrente; l'MPC le calcola con questo.
  \item \emph{Si considera una cella sola.} Le celle sono indipendenti, quindi
        l'indice $n$ è sottinteso: ogni formula va applicata a ciascuna cella,
        con i suoi parametri.
  \item \emph{La corrente è un parametro.} Come visto nella
        sezione~\ref{sec:modello}, sull'orizzonte $I$ è costante. Le derivate
        della soluzione di un'equazione differenziale rispetto a una costante
        da cui essa dipende si chiamano \emph{derivate parametriche}, o
        sensibilità.
\end{itemize}
La difficoltà è che la tensione ai morsetti dipende dalla corrente in due modi:
direttamente, attraverso la caduta $R_o I$, e attraverso lo stato $(s, V_p)$,
che la corrente ha spostato lungo l'orizzonte. Per derivarla bisogna quindi
sapere prima di quanto la corrente sposta lo stato. Il calcolo procede in tre
tempi:
\begin{itemize}
  \item la sezione~\ref{sec:sens} introduce le derivate parametriche dello
        stato, $\alpha$ e $\beta$, e ricava a tempo continuo le equazioni a cui
        obbediscono;
  \item la sezione~\ref{sec:uso} mostra, sempre a tempo continuo, come da
        $\alpha$ e $\beta$ si arriva alla retta della \eqref{eq:lin} e quindi
        al QP;
  \item la sezione~\ref{sec:derdiscreto} discretizza il tutto e arriva alle
        formule che il codice applica.
\end{itemize}

\subsection{A tempo continuo: equazioni di sensibilità}\label{sec:sens}

\subsubsection{Le derivate parametriche dello stato}

Con la corrente $I$ costante, la soluzione della \eqref{eq:c-compatta} è una
funzione del tempo e del valore scelto per la corrente: $\zeta(t;\, I)$. Le
derivate parametriche dello stato sono
\[
  \alpha(t) = \frac{\partial s(t;\, I)}{\partial I}, \qquad
  \beta(t)  = \frac{\partial V_p(t;\, I)}{\partial I} .
\]
Dicono di quanto si sposta lo stato all'istante $t$ se la corrente, costante su
tutto l'orizzonte, aumenta di un ampere. Insieme alla derivata della
tensione ai morsetti, $d$, che si ricava da esse, sono i tre
\emph{coefficienti di sensibilità} della cella: $\alpha$ del SOC, $\beta$
della tensione di polarizzazione, $d$ della tensione ai morsetti. Tutti e tre
sono derivate parametriche rispetto alla corrente. Nel seguito sono usati
così:
\begin{itemize}
  \item $\alpha$ è già uno dei due coefficienti che entrano nella
        \eqref{eq:lin}: negli istanti di campionamento è $\alpha_{j,n}$, la
        derivata del SOC (\texttt{alfa} nel codice);
  \item $\alpha$ e $\beta$ insieme danno la derivata della tensione, cioè $d$
        (\texttt{d} nel codice): con la \eqref{eq:c-dy} a tempo continuo e
        con la \eqref{eq:D} a tempo discreto;
  \item $\alpha$ serve anche a calcolare $\beta$, perché i parametri del ramo
        RC dipendono dal SOC: compare nella \eqref{eq:c-beta} e nella
        \eqref{eq:beta};
  \item $d$ e $\alpha$ entrano infine nel costo e nei vincoli del QP
        (sezione~\ref{sec:qp}).
\end{itemize}
$\beta$ non esce quindi dal calcolo della predizione: è un passaggio intermedio
verso $d$. Nel codice, al passo $j$, $\alpha_j$, $\beta_j$ e $d_j$ sono
\file{alfa_j}, \file{beta_j} e \file{d_j}.

Va tenuto presente che $\alpha$ e $\beta$ non sono derivate rispetto al tempo.
Sono a loro volta funzioni del tempo: partono da zero e cambiano lungo
l'orizzonte, man mano che l'effetto della corrente si accumula. Come lo stato,
obbediscono a una propria equazione differenziale, detta equazione di
sensibilità. Ha la stessa forma per qualunque sistema, quindi conviene
scriverla prima in generale.

\subsubsection{Espressione generale}

Sia dato un sistema qualsiasi, con stato $x \in \mathbb{R}^{n_x}$ e un vettore
$\theta \in \mathbb{R}^{n_\theta}$ di parametri costanti nel tempo:
\[
  \dot x = f(x, \theta), \qquad y = g(x, \theta), \qquad x(0) = x_0(\theta) .
\]
La matrice $n_x \times n_\theta$ delle derivate parametriche dello stato,
$\partial x(t;\theta)/\partial\theta$, è la soluzione dell'equazione
differenziale
\[
  \frac{d}{dt}\,\frac{\partial x}{\partial \theta}
    = \frac{\partial f}{\partial x}\,\frac{\partial x}{\partial \theta}
      + \frac{\partial f}{\partial \theta} ,
  \qquad
  \frac{\partial x}{\partial \theta}(0) = \frac{\partial x_0}{\partial \theta} ,
\]
e da essa si ricava, senza altre integrazioni, la derivata parametrica
dell'uscita:
\[
  \frac{\partial y}{\partial \theta}
    = \frac{\partial g}{\partial x}\,\frac{\partial x}{\partial \theta}
      + \frac{\partial g}{\partial \theta} .
\]
Le due formule si ottengono derivando rispetto a $\theta$ le equazioni del
sistema e, nella prima, scambiando l'ordine delle derivate rispetto a $t$ e a
$\theta$. In entrambe il parametro entra in due modi: attraverso lo stato, che
esso stesso ha spostato (primo addendo), e direttamente (secondo addendo). Tre
proprietà dell'equazione di sensibilità servono nel seguito.
\begin{itemize}
  \item È lineare nell'incognita $\partial x/\partial\theta$, anche quando il
        sistema non lo è.
  \item I suoi coefficienti sono le jacobiane di $f$, valutate lungo la
        traiettoria $x(t;\theta)$: cambiano nel tempo, e per conoscerle serve
        lo stato.
  \item La dipendenza va in una sola direzione: le derivate hanno bisogno dello
        stato, lo stato non ha bisogno delle derivate. Lo stato si può quindi
        calcolare da solo; le derivate si calcolano insieme allo stato, oppure
        dopo.
\end{itemize}
Un esempio eseguibile di queste formule, su un sistema massa-molla-smorzatore
integrato con \texttt{ode45}, è in \file{esempio_sensibilita_ode45.m}.

\subsubsection{Applicazione alla cella}

Per la cella il parametro è uno solo, $\theta = I$; lo stato è $\zeta = (s, V_p)$,
quindi $\partial \zeta/\partial I = (\alpha, \beta)$; la condizione iniziale $\zeta_k$
è misurata, quindi non dipende da $I$. L'espressione generale diventa
\begin{equation}
  \frac{d}{dt}\,\frac{\partial \zeta}{\partial I}
    = \underbrace{\frac{\partial f}{\partial \zeta}\,
                  \frac{\partial \zeta}{\partial I}}_{I \text{ tramite lo stato}}
    \;+\; \underbrace{\frac{\partial f}{\partial I}}_{I \text{ diretta}} ,
  \qquad \frac{\partial \zeta}{\partial I}(0) = 0 .
  \tag{C5}\label{eq:c-sens}
\end{equation}
La condizione iniziale è nulla perché lo stato misurato non dipende dalla
corrente futura. Scritta per esteso, la \eqref{eq:c-sens} è
\[
  \frac{d}{dt}
  \begin{bmatrix} \alpha \\[4pt] \beta \end{bmatrix}
  =
  \begin{bmatrix}
    \dfrac{\partial f_1}{\partial s} & \dfrac{\partial f_1}{\partial V_p} \\[10pt]
    \dfrac{\partial f_2}{\partial s} & \dfrac{\partial f_2}{\partial V_p}
  \end{bmatrix}
  \begin{bmatrix} \alpha \\[4pt] \beta \end{bmatrix}
  +
  \begin{bmatrix}
    \dfrac{\partial f_1}{\partial I} \\[10pt] \dfrac{\partial f_2}{\partial I}
  \end{bmatrix},
\]
dove $f_1$ e $f_2$ sono i secondi membri della \eqref{eq:c-soc} e della
\eqref{eq:c-vp}:
\[
  f_1 = -\frac{\eta}{3600\,C}\, I , \qquad
  f_2 = -\frac{V_p}{\tau(s)} + \frac{I}{C_p(s)} .
\]
Ogni riga è l'equazione di una delle due derivate: la prima di $\alpha$, la
seconda di $\beta$. Restano da calcolare le sei derivate parziali.
\begin{itemize}
  \item $f_1$ contiene solo la corrente, quindi
        $\partial f_1/\partial s = \partial f_1/\partial V_p = 0$ e
        $\partial f_1/\partial I = -\eta/(3600\,C)$: la velocità con cui la
        cella si scarica dipende dalla corrente, non dallo stato.
  \item In $f_2$ la tensione $V_p$ compare solo al numeratore del primo
        termine e la corrente solo al numeratore del secondo, quindi
        $\partial f_2/\partial V_p = -1/\tau$ e
        $\partial f_2/\partial I = 1/C_p$.
  \item Il SOC compare in $f_2$ solo dentro $\tau$ e $C_p$, che stanno al
        denominatore. Poiché la derivata di $1/\tau(s)$ rispetto a $s$ è
        $-\tau'/\tau^2$, e lo stesso vale per $1/C_p(s)$,
        \[
          \frac{\partial f_2}{\partial s}
            = \frac{V_p\,\tau'}{\tau^2} - \frac{I\,C_p'}{C_p^2} .
        \]
\end{itemize}
In forma di matrici:
\[
  \frac{\partial f}{\partial \zeta} =
  \begin{bmatrix}
    0 & 0 \\[4pt]
    \dfrac{V_p\,\tau'}{\tau^2} - \dfrac{I\,C_p'}{C_p^2} & -\dfrac{1}{\tau}
  \end{bmatrix},
  \qquad
  \frac{\partial f}{\partial I} =
  \begin{bmatrix}
    -\dfrac{\eta}{3600\,C} \\[8pt] \dfrac{1}{C_p}
  \end{bmatrix}.
\]
In queste formule, e in quelle che seguono, l'argomento è sottinteso per non
appesantire la scrittura: $\tau$, $C_p$, $\tau'$ e $C_p'$ stanno per
$\tau(s(t))$, $C_p(s(t))$ e le loro derivate rispetto a $s$, e $V_p$ sta per
$V_p(t)$. Sono grandezze che cambiano nel tempo, perché lo stato si muove lungo
la traiettoria; la corrente $I$ invece è costante. A ogni istante sono comunque
numeri noti: lo stato viene dalle \eqref{eq:c-soc}--\eqref{eq:c-vp}, i
parametri e le loro derivate dalle tabelle con la \eqref{eq:tabella}, e
$\tau' = R_p'\,C_p + R_p\,C_p'$ dalla regola di derivazione del prodotto
$\tau = R_p\,C_p$.

Sostituendo, la prima riga della \eqref{eq:c-sens} dà la \eqref{eq:c-alpha} e
la seconda la \eqref{eq:c-beta}:
\begin{align}
  \dot\alpha(t) &= -\frac{\eta}{3600\,C} ,
     \tag{C6}\label{eq:c-alpha}\\
  \dot\beta(t) &=
      \underbrace{-\frac{\beta}{\tau}}_{V_p \text{ dipende da } I}
      \;+\; \underbrace{\frac{V_p\,\tau'\,\alpha}{\tau^2}}_{\tau \text{ tramite il SOC}}
      \;+\; \underbrace{\frac{1}{C_p}}_{I \text{ diretta}}
      \;-\; \underbrace{\frac{I\,C_p'\,\alpha}{C_p^2}}_{C_p \text{ tramite il SOC}} ,
     \tag{C7}\label{eq:c-beta}
\end{align}
con $\alpha(0) = \beta(0) = 0$. Per l'uscita non serve un'altra equazione
differenziale: basta la seconda formula generale. Con
$g = V_{oc}(s) - V_p - R_o(s)\,I$ le derivate parziali sono
$\partial g/\partial s = V_{oc}' - R_o'\,I$,
$\partial g/\partial V_p = -1$ e $\partial g/\partial I = -R_o$, quindi
\begin{equation}
  \frac{\partial y(t)}{\partial I}
    = V_{oc}'\,\alpha \;-\; \beta \;-\; R_o'\,\alpha\,I \;-\; R_o .
  \tag{C8}\label{eq:c-dy}
\end{equation}
Tre osservazioni.
\begin{itemize}
  \item La \eqref{eq:c-alpha} si integra subito: $\alpha(t) = -\eta\,t/(3600\,C)$.
        Il SOC è esattamente lineare nella corrente. Fa eccezione la cella
        satura: finché resta a $s = 1$ lo fa per qualunque corrente vicina,
        quindi $\alpha = 0$.
  \item La \eqref{eq:c-beta} è lineare in $\beta$, ma i suoi coefficienti
        contengono $s(t)$, $V_p(t)$ e $\alpha(t)$: non si può risolvere da
        sola, va integrata dopo lo stato o insieme a esso, lungo la stessa
        traiettoria. Lo stato invece non contiene né $\alpha$ né $\beta$: è la
        dipendenza in una sola direzione già vista in generale.
  \item Se i parametri non dipendessero dal SOC ($\tau' = C_p' = 0$)
        resterebbe $\dot\beta = -\beta/\tau + 1/C_p$, la risposta di un
        circuito RC lineare. I due termini con $\alpha$ sono il contributo
        della non linearità.
\end{itemize}

\subsection{A tempo continuo: uso di $\alpha$ e $\beta$}\label{sec:uso}

Le \eqref{eq:c-alpha}--\eqref{eq:c-dy} si calcolano lungo una traiettoria
precisa: la traiettoria di linearizzazione della sezione~\ref{sec:punto}, cioè
quella che la cella percorre partendo dallo stato misurato, con la corrente
$I = i_k + \bar u$ tenuta costante sull'orizzonte. A tempo continuo la si
indica con $\bar s(t)$, $\bar V_p(t)$, $\bar y(t)$. Lungo questa traiettoria
$\alpha(t)$ e $\beta(t)$ sono funzioni note del tempo: si calcolano una volta
sola, prima di ottimizzare.

L'MPC deve però confrontare correnti diverse: $i_k + u_k$, con $u_k$ ancora da
scegliere. Ognuna differisce dalla corrente di linearizzazione di
$u_k - \bar u$. È qui che servono $\alpha$ e $\beta$.

\paragraph{Primo uso: lo stato.} $\alpha$ e $\beta$ dicono dove si troverebbe
lo stato con una di queste correnti, senza risolvere di nuovo le equazioni
differenziali:
\begin{align}
  s(t;\, u_k) &\simeq \bar s(t) + \alpha(t)\,\bigl(u_k - \bar u\bigr) ,
     \tag{C9}\label{eq:c-s-lin}\\
  V_p(t;\, u_k) &\simeq \bar V_p(t) + \beta(t)\,\bigl(u_k - \bar u\bigr) .
     \tag{C10}\label{eq:c-vp-lin}
\end{align}
Sono sviluppi al primo ordine: lo spostamento dello stato è la derivata
moltiplicata per la variazione di corrente. Per il SOC la \eqref{eq:c-s-lin} è
esatta finché la cella non satura; per $V_p$ la \eqref{eq:c-vp-lin} è
un'approssimazione, buona se $u_k$ è vicina a $\bar u$.

\paragraph{Secondo uso: l'uscita.} È quello che interessa all'MPC. Sia $d(t)$
la derivata dell'uscita \eqref{eq:c-dy}, calcolata lungo la traiettoria di
linearizzazione. Allora
\begin{equation}
  y(t;\, u_k) \;\simeq\; \bar y(t) + d(t)\,\bigl(u_k - \bar u\bigr) ,
  \qquad
  d(t) = V_{oc}'\,\alpha(t) - \beta(t) - R_o'\,\alpha(t)\,I - R_o .
  \tag{C11}\label{eq:c-y-lin}
\end{equation}
È la versione a tempo continuo della prima riga della \eqref{eq:lin}; la
\eqref{eq:c-s-lin} lo è della seconda. Per ogni istante $t$ la
\eqref{eq:c-y-lin} è una retta in $u_k$: la tangente, nel punto $\bar u$, alla
curva vera $y(t;\, u_k)$. I due coefficienti $\bar y(t)$ e $d(t)$ non dipendono
da $u_k$; l'incognita compare solo nel fattore $(u_k - \bar u)$. In $d(t)$ ogni
termine ha un ruolo:
\begin{itemize}
  \item $V_{oc}'\,\alpha$ e $R_o'\,\alpha\,I$ dicono di quanto cambiano la
        tensione a vuoto e la caduta sulla resistenza serie perché la corrente
        ha cambiato il SOC;
  \item $\beta$ dice di quanto è cambiata la caduta sul ramo RC;
  \item $R_o$ è l'effetto immediato della corrente sulla resistenza serie. È
        il solo presente all'istante iniziale: $d(0) = -R_o$, perché
        $\alpha(0) = \beta(0) = 0$. Gli altri crescono lungo l'orizzonte.
\end{itemize}

\paragraph{Verso il QP.} L'MPC sostituisce la predizione vera con queste
rette. Per vedere che cosa ne guadagna basta scrivere, a tempo continuo, un
vincolo di tensione minima e un costo che fa inseguire a ogni cella la
tensione $y^0(t)$ della cella nominale:
\begin{gather*}
  \bar y^n(t) + d^n(t)\,\bigl(u^n_k - \bar u^n\bigr) \;\ge\; y_{\min}
  \qquad \text{per ogni } t \in [0,\; p\,\Tc] , \\[4pt]
  \int_0^{p\,\Tc} \sum_{n=1}^{N}
     \Bigl(\bar y^n(t) + d^n(t)\,\bigl(u^n_k - \bar u^n\bigr) - y^0(t)\Bigr)^2 dt .
\end{gather*}
Il vincolo è lineare in $u_k$ e il costo quadratico: è già la struttura di un
QP. In questa forma però i vincoli sono infiniti, uno per ogni istante, e il
costo è un integrale. Il problema dell'MPC (eq.~6) è invece formulato solo agli
istanti di campionamento $k+1, \dots, k+p$: bastano i valori di $\bar y$, $d$
e $\alpha$ in quegli istanti. Li fornisce la discretizzazione.

\subsection{Discretizzazione temporale}\label{sec:derdiscreto}

La sezione precedente si è chiusa con una richiesta precisa: servono
$\bar y$, $d$ e $\alpha$ solo negli istanti di campionamento. Qui li si
calcola, arrivando alle formule che il codice applica: le \eqref{eq:beta},
\eqref{eq:alpha} e \eqref{eq:D}.

\subsubsection{Istanti di campionamento e grandezze discrete}

L'orizzonte $[0,\; p\,\Tc]$ è diviso in $p$ passi di durata $\Tc$. Gli istanti
di campionamento sono $t_j = j\,\Tc$, con $j = 0, \dots, p$: $t_0$ è l'istante
attuale $k$ e $t_j$ l'istante $k+j$. Il passo $j$ è l'intervallo tra
$t_{j-1}$ e $t_j$: comincia all'istante $k+j-1$ e finisce all'istante $k+j$.

Le derivate parametriche dello stato negli istanti di campionamento sono
\[
  \alpha_j = \frac{\partial s_{k+j}}{\partial I}, \qquad
  \beta_j  = \frac{\partial V_{p,k+j}}{\partial I}, \qquad
  \alpha_0 = \beta_0 = 0 ,
\]
con condizioni iniziali nulle perché lo stato misurato $(s_k, V_{p,k})$ non
dipende dalla corrente futura. A ogni grandezza a tempo continuo corrisponde
così una grandezza discreta, con lo stesso significato:

\begin{center}
\begin{tabular}{llll}
\toprule
 & Tempo continuo, in $t_j$ & Tempo discreto & Nel codice \\
\midrule
Stato & $\bar s(t_j)$, $\bar V_p(t_j)$ & $s_{k+j}$, $V_{p,k+j}$ & \texttt{s}, \texttt{Vp} \\
Uscita & $\bar y(t_j)$ & $y_{k+j}$ & \texttt{y} \\
Sensibilità di $s$ e $V_p$ & $\alpha(t_j)$, $\beta(t_j)$ & $\alpha_j$, $\beta_j$ & \file{alfa_j}, \file{beta_j} \\
Sensibilità di $y$ & $d(t_j)$ & $d_j$ & \file{d_j} \\
\bottomrule
\end{tabular}
\end{center}

\noindent
Le grandezze discrete non sono i valori esatti di quelle continue negli istanti
$t_j$, ma la loro approssimazione con il metodo di Eulero; solo per il SOC e
per $\alpha$ le due coincidono. Raccolte per tutti i passi e per tutte le
celle, $y_{k+j}$ e $s_{k+j}$ formano le matrici $\bar Y$ e $\bar S$ della
sezione~\ref{sec:punto} (nel codice \file{Y_bar} e \file{S_bar}), $d_j$ e
$\alpha_j$ i coefficienti $d_{j,n}$ e $\alpha_{j,n}$ della \eqref{eq:lin}
(matrici \texttt{d} e \texttt{alfa}).

\subsubsection{Due strade equivalenti}

Per passare dalle equazioni di sensibilità a formule discrete ci sono due
strade:
\begin{enumerate}
  \item discretizzare con Eulero in avanti le equazioni di sensibilità
        \eqref{eq:c-alpha}--\eqref{eq:c-beta}, con lo stesso passo $\Tc$ usato
        per lo stato;
  \item derivare rispetto a $I$ le equazioni dello stato già discretizzate,
        \eqref{eq:soc}--\eqref{eq:y}.
\end{enumerate}
La regola di Eulero è quella della sezione~\ref{sec:discreto}: una grandezza
$\xi(t)$ avanza di un passo come se la sua derivata, calcolata all'\emph{inizio}
del passo, restasse costante per tutta la durata del passo,
\[
  \xi_j = \xi_{j-1} + \Tc\; \dot \xi\,\big|_{t_{j-1}} .
\]
Con Eulero in avanti le due strade portano alle stesse formule. Derivando un
passo del modello discreto si ottiene infatti
\[
  \zeta_{k+j} = \zeta_{k+j-1} + \Tc\, f(\zeta_{k+j-1}, I)
  \quad\Longrightarrow\quad
  \frac{\partial \zeta_{k+j}}{\partial I}
    = \frac{\partial \zeta_{k+j-1}}{\partial I}
      + \Tc\left(\frac{\partial f}{\partial \zeta}\,
                 \frac{\partial \zeta_{k+j-1}}{\partial I}
                 + \frac{\partial f}{\partial I}\right),
\]
che è proprio un passo di Eulero sulla \eqref{eq:c-sens}.

Il codice segue la seconda strada: a ogni passo, accanto a ogni equazione del
modello c'è la sua derivata, e le derivate avanzano insieme allo stato. La
conseguenza è importante: le derivate che si ottengono sono quelle
\emph{esatte} del modello discreto, cioè proprio del modello su cui l'MPC
costruisce costo e vincoli. Non c'è un secondo errore di approssimazione che
si aggiunge a quello del modello.

Di seguito le tre equazioni del modello discreto sono derivate una alla volta,
nell'ordine in cui il codice le applica: prima la tensione di polarizzazione,
poi lo stato di carica, infine la tensione ai morsetti. Per ciascuna si segue
la seconda strada e si verifica che la prima dà lo stesso risultato.

\subsubsection{Tensione di polarizzazione}

Si cerca $\beta_j$, la derivata rispetto a $I$ della tensione di polarizzazione
a fine passo, $V_{p,k+j}$. Il punto di partenza è l'equazione che la calcola,
la \eqref{eq:vp}:
\[
  V_{p,k+j} = V_{p,k+j-1} + \Tc\left(-\frac{V_{p,k+j-1}}{\tau(s_{k+j-1})}
              + \frac{I}{C_p(s_{k+j-1})}\right) .
\]
A secondo membro tutto è riferito all'\emph{inizio} del passo: la tensione
$V_{p,k+j-1}$ e i parametri del ramo RC, letti nel SOC $s_{k+j-1}$.

\paragraph{Notazione.} Poiché in tutto questo paragrafo i parametri sono
valutati nello stesso punto, $s_{k+j-1}$, l'argomento viene sottinteso:
\[
  \tau = \tau(s_{k+j-1}), \qquad C_p = C_p(s_{k+j-1}), \qquad
  \tau' = \frac{d\tau}{ds}(s_{k+j-1}), \qquad
  C_p' = \frac{dC_p}{ds}(s_{k+j-1}) .
\]
Con questa scrittura la \eqref{eq:vp} diventa
$V_{p,k+j} = V_{p,k+j-1} + \Tc\,\bigl(-V_{p,k+j-1}/\tau + I/C_p\bigr)$.

\paragraph{Che cosa dipende dalla corrente.} Prima di derivare bisogna
stabilire quali grandezze, a secondo membro, cambiano se cambia $I$. Sono tre.
\begin{itemize}
  \item La corrente $I$ stessa, che compare in modo esplicito al numeratore
        dell'ultimo termine.
  \item La tensione di inizio passo, $V_{p,k+j-1}$: è il risultato dei passi
        precedenti, durante i quali la corrente ha già agito. La sua derivata
        è, per definizione, $\beta_{j-1}$.
  \item I parametri $\tau$ e $C_p$: non contengono $I$ in modo esplicito, ma
        sono letti nel SOC $s_{k+j-1}$, che la corrente ha spostato nei passi
        precedenti. La derivata di $s_{k+j-1}$ è, per definizione,
        $\alpha_{j-1}$; per la regola di derivazione delle funzioni composte
        \[
          \frac{\partial \tau}{\partial I}
             = \frac{d\tau}{ds}\,\frac{\partial s_{k+j-1}}{\partial I}
             = \tau'\,\alpha_{j-1} ,
          \qquad
          \frac{\partial C_p}{\partial I}
             = \frac{dC_p}{ds}\,\frac{\partial s_{k+j-1}}{\partial I}
             = C_p'\,\alpha_{j-1} .
        \]
\end{itemize}
Il passo $\Tc$ invece è una costante. Il terzo punto è quello in cui la non
linearità del modello entra nel calcolo: se i parametri non dipendessero dal
SOC, $\tau'$ e $C_p'$ sarebbero nulli e la corrente agirebbe solo nei primi due
modi.

\paragraph{Derivata termine per termine.} I tre addendi della \eqref{eq:vp} si
derivano uno alla volta. Il secondo e il terzo sono quozienti, e sia il
numeratore sia il denominatore possono dipendere da $I$:
\begin{align*}
  \frac{\partial}{\partial I}\, V_{p,k+j-1} &= \beta_{j-1} , \\[4pt]
  \frac{\partial}{\partial I} \left(-\frac{V_{p,k+j-1}}{\tau}\right)
     &= -\frac{1}{\tau}\,\frac{\partial V_{p,k+j-1}}{\partial I}
        + \frac{V_{p,k+j-1}}{\tau^2}\,\frac{\partial \tau}{\partial I}
      = -\frac{\beta_{j-1}}{\tau}
        + \frac{V_{p,k+j-1}\,\tau'\,\alpha_{j-1}}{\tau^2} , \\[4pt]
  \frac{\partial}{\partial I} \left(\frac{I}{C_p}\right)
     &= \frac{1}{C_p} - \frac{I}{C_p^2}\,\frac{\partial C_p}{\partial I}
      = \frac{1}{C_p} - \frac{I\,C_p'\,\alpha_{j-1}}{C_p^2} .
\end{align*}
Nella seconda e nella terza riga il primo addendo viene dal numeratore, il
secondo dal denominatore.

\paragraph{Risultato.} Sommando le tre righe, con il fattore $\Tc$ davanti
alle ultime due:
\begin{equation}
  \beta_j = \beta_{j-1} + \Tc\left(
      \underbrace{-\frac{\beta_{j-1}}{\tau}}_{V_p \text{ dipende da } I}
      \;+\; \underbrace{\frac{V_{p,k+j-1}\,\tau'\,\alpha_{j-1}}{\tau^2}}_{\tau \text{ tramite il SOC}}
      \;+\; \underbrace{\frac{1}{C_p}}_{I \text{ diretta}}
      \;-\; \underbrace{\frac{I\,C_p'\,\alpha_{j-1}}{C_p^2}}_{C_p \text{ tramite il SOC}}
  \right) .
  \label{eq:beta}
\end{equation}
I quattro termini tra parentesi sono i modi in cui la corrente agisce su
$V_{p,k+j}$ durante il passo: attraverso la tensione di inizio passo,
attraverso $\tau$, in modo diretto e attraverso $C_p$. Nel codice è
l'istruzione
{\small
\begin{verbatim}
    beta_j = beta_j + Tc*( -beta_j/tau + Vp*ds_tau*alfa_j/tau^2 ...
                           + 1/Cp - I_cella*ds_Cp*alfa_j/Cp^2 );
\end{verbatim}
}
\noindent
Al primo passo, $j = 1$, le derivate di inizio passo sono $\alpha_0 = \beta_0
= 0$: tre termini su quattro si annullano e resta
$\beta_1 = \Tc / C_p(s_k)$. Dopo un solo passo la corrente ha agito sulla
tensione di polarizzazione solo in modo diretto; gli altri tre termini
compaiono dal secondo passo in poi.

\paragraph{Verifica con la prima strada.} La stessa formula si ottiene
discretizzando la \eqref{eq:c-beta} con Eulero: si sostituisce $\dot\beta$ con
il rapporto incrementale e si valuta il secondo membro all'inizio del passo,
cioè in $t_{j-1}$:
\[
  \frac{\beta_j - \beta_{j-1}}{\Tc}
    = -\frac{\beta_{j-1}}{\tau}
      + \frac{V_{p,k+j-1}\,\tau'\,\alpha_{j-1}}{\tau^2}
      + \frac{1}{C_p}
      - \frac{I\,C_p'\,\alpha_{j-1}}{C_p^2} .
\]
Moltiplicando per $\Tc$ e portando $\beta_{j-1}$ a secondo membro si ritrova
la \eqref{eq:beta}, come previsto.

\paragraph{Dai simboli ai numeri.} Quando il codice applica la
\eqref{eq:beta} al passo $j$, ogni simbolo a secondo membro è un numero già
disponibile:
\begin{itemize}
  \item $\Tc$ e $I$ sono fissati prima di cominciare: il passo di controllo e
        la corrente del punto di linearizzazione;
  \item $V_{p,k+j-1}$, $\alpha_{j-1}$ e $\beta_{j-1}$ sono stati calcolati al
        passo precedente; al primo passo sono la tensione misurata $V_{p,k}$ e
        due zeri;
  \item $R_p$, $C_p$ e le loro derivate $R_p'$, $C_p'$ si leggono dalle
        tabelle con la \eqref{eq:tabella}, nel SOC $s_{k+j-1}$;
  \item da questi si calcolano $\tau$ e la sua derivata, con la regola di
        derivazione di un prodotto:
        $\tau = R_p\,C_p$ e $\tau' = R_p'\,C_p + R_p\,C_p'$.
\end{itemize}
In nessun punto serve una derivata calcolata per via numerica, per esempio con
le differenze finite. L'esempio numerico più avanti ripercorre questi passaggi
con i valori del progetto.

\subsubsection{Stato di carica}

Si cerca $\alpha_j$, la derivata rispetto a $I$ del SOC a fine passo,
$s_{k+j}$. L'equazione di partenza è la \eqref{eq:soc},
$s_{k+j} = \min\bigl(s_{k+j-1} - \kappa I,\; 1\bigr)$, in cui $\kappa$ è una
costante della cella: non dipende né dalla corrente né dal SOC. Ci sono due
casi.
\begin{itemize}
  \item \emph{La cella non satura}: $s_{k+j} = s_{k+j-1} - \kappa I$. La
        corrente agisce attraverso il SOC di inizio passo, la cui derivata è
        $\alpha_{j-1}$, e in modo diretto, con derivata $-\kappa$.
  \item \emph{La cella satura}: il minimo sceglie la costante 1. Il SOC resta
        a 1 anche per correnti un poco diverse, quindi non dipende più dalla
        corrente e la sua derivata è nulla.
\end{itemize}
In tutto:
\begin{equation}
  \alpha_j =
  \begin{cases}
    \alpha_{j-1} - \kappa & \text{se } s_{k+j-1} - \kappa I < 1,\\[2pt]
    0 & \text{se la cella satura a } s_{k+j} = 1 .
  \end{cases}
  \label{eq:alpha}
\end{equation}
La prima strada dà lo stesso risultato: Eulero sulla \eqref{eq:c-alpha} è
$\alpha_j = \alpha_{j-1} + \Tc \cdot \bigl(-\eta/(3600\,C)\bigr)
= \alpha_{j-1} - \kappa$. Finché la cella non satura, quindi,
$\alpha_j = -j\,\kappa$, che coincide con la soluzione esatta della
\eqref{eq:c-alpha} all'istante $t_j$: qui Eulero non introduce errore, perché
il secondo membro della \eqref{eq:c-alpha} è costante. È la conferma, a tempo
discreto, che il SOC è esattamente lineare nella corrente.

\subsubsection{Tensione ai morsetti}

Si cercano i due coefficienti di sensibilità del passo $j$ che entrano nella
\eqref{eq:lin}. Quello del SOC è $\alpha_j$, appena calcolato con la
\eqref{eq:alpha}. Per quello della tensione, $d_j$, il punto di partenza è
l'equazione dell'uscita, la \eqref{eq:y}:
\[
  y_{k+j} = V_{oc}(s_{k+j}) - V_{p,k+j} - R_o(s_{k+j})\, I .
\]
A differenza della \eqref{eq:vp}, qui tutto è riferito alla \emph{fine} del
passo, perché $y_{k+j}$ è la tensione all'istante $k+j$: compaiono lo stato
appena aggiornato e i parametri letti nel nuovo SOC. Di conseguenza anche le
derivate che servono sono quelle di fine passo, $\alpha_j$ e $\beta_j$.
Sottintendendo l'argomento come prima, $V_{oc}$, $R_o$, $V_{oc}'$ e $R_o'$ sono
valutati in $s_{k+j}$.

La corrente agisce in quattro modi:
\begin{itemize}
  \item attraverso $V_{oc}$, che dipende dal SOC:
        $\partial V_{oc}/\partial I = V_{oc}'\,\alpha_j$;
  \item attraverso $V_{p,k+j}$, la cui derivata è $\beta_j$;
  \item attraverso $R_o$, che dipende dal SOC:
        $\partial R_o/\partial I = R_o'\,\alpha_j$;
  \item in modo diretto, perché $I$ moltiplica $R_o$.
\end{itemize}
Gli ultimi due vengono dalla derivata del prodotto $R_o\,I$, che è
$R_o'\,\alpha_j\,I + R_o$. Sommando:
\begin{equation}
\begin{aligned}
  d_j &= \frac{\partial y_{k+j}}{\partial I}
      = \underbrace{V_{oc}'\,\alpha_j}_{V_{oc} \text{ tramite il SOC}}
        \;-\; \underbrace{\beta_j}_{V_p \text{ dipende da } I}
        \;-\; \underbrace{R_o'\,\alpha_j\,I}_{R_o \text{ tramite il SOC}}
        \;-\; \underbrace{R_o}_{I \text{ diretta}} .
\end{aligned}
\label{eq:D}
\end{equation}
Nel codice è l'istruzione
{\small
\begin{verbatim}
    d_j = ds_Voc*alfa_j - beta_j - ds_Ro*alfa_j*I_cella - Ro;
\end{verbatim}
}
\noindent
L'uscita non ha una dinamica propria, quindi qui non c'è nulla da discretizzare: la
\eqref{eq:D} è la \eqref{eq:c-dy} calcolata nello stato di fine passo, ed è il
valore discreto di $d(t_j)$. I numeri si ottengono come prima: $\alpha_j$ e
$\beta_j$ sono appena stati aggiornati, mentre $V_{oc}$, $R_o$ e le loro
derivate si leggono dalle tabelle con la \eqref{eq:tabella}, questa volta in
$s_{k+j}$.

\subsubsection{Ordine dei calcoli}

Mettendo insieme le tre derivazioni, un passo $j$ dell'orizzonte si svolge
così:
\begin{enumerate}
  \item si leggono dalle tabelle, con la \eqref{eq:tabella}, $R_p$, $C_p$ e le
        loro derivate nel SOC di inizio passo, $s_{k+j-1}$, e da questi si
        calcolano $\tau$ e $\tau'$;
  \item si aggiorna $\beta$ con la \eqref{eq:beta}, che usa $V_{p,k+j-1}$,
        $\alpha_{j-1}$ e $\beta_{j-1}$;
  \item si aggiorna $V_p$ con la \eqref{eq:vp};
  \item si aggiornano $s$ con la \eqref{eq:soc} e $\alpha$ con la
        \eqref{eq:alpha};
  \item si leggono allo stesso modo $V_{oc}$, $R_o$ e le loro derivate nel
        nuovo SOC, $s_{k+j}$;
  \item si calcolano $y_{k+j}$ con la \eqref{eq:y} e $d_j$ con la
        \eqref{eq:D}; $y_{k+j}$, $s_{k+j}$, $d_j$ e $\alpha_j$ si scrivono
        nella riga $j$ delle quattro matrici.
\end{enumerate}
L'ordine non è arbitrario. Nel codice ogni grandezza occupa una sola
variabile, che a ogni passo viene sovrascritta: prima dell'aggiornamento
contiene il valore di inizio passo, dopo quello di fine passo. La
\eqref{eq:beta} vuole $V_p$ e $\alpha$ di \emph{inizio} passo, quindi $\beta$
va aggiornata prima di loro; la \eqref{eq:D} vuole stato e derivate di
\emph{fine} passo, quindi viene per ultima. Al primo passo si parte dallo
stato misurato e da $\alpha_0 = \beta_0 = 0$; a ogni passo successivo dai
valori lasciati dal precedente.

\subsubsection{Esempio numerico}

Si considera una cella nominale, cioè con tutti i fattori di deviazione uguali
a 1, e i valori di \file{init_parametri.m}: $\Tc = 20$~s e
$\kappa = 2{,}011 \cdot 10^{-4}~\mathrm{A}^{-1}$. Lo stato misurato è
$s_k = 0{,}6$, $V_{p,k} = 0$; la corrente del punto di linearizzazione è
$I = 10$~A.

\paragraph{Primo passo.} Le tabelle si leggono in $s_k = 0{,}6$, nel tratto
tra $0{,}5$ e $0{,}75$. Per $R_p$ il calcolo è quello della
sezione~\ref{sec:tabelle}: $R_p = 0{,}00188~\Omega$ e
$R_p' = 0{,}0028~\Omega$. Allo stesso modo, per $C_p$:
\[
  C_p' = \frac{33630 - 18721}{0{,}75 - 0{,}5} = 59636~\mathrm{F} ,
  \qquad
  C_p = 18721 + 59636 \cdot 0{,}1 = 24685~\mathrm{F} .
\]
Da questi quattro numeri:
\[
  \tau = R_p\,C_p = 46{,}4~\mathrm{s} ,
  \qquad
  \tau' = R_p'\,C_p + R_p\,C_p'
        = 0{,}0028 \cdot 24685 + 0{,}00188 \cdot 59636
        = 181{,}2~\mathrm{s} .
\]
Come per $R_p'$, la derivata è riferita all'intera escursione del SOC:
$\tau' = 181{,}2$~s significa che, se il SOC scende di $0{,}01$, $\tau$ scende
di circa $1{,}8$~s. Poiché $\alpha_0 = \beta_0 = 0$, della \eqref{eq:beta}
resta solo il termine diretto, e lo stato avanza con le \eqref{eq:vp} e
\eqref{eq:soc}:
\begin{align*}
  \beta_1 &= \frac{\Tc}{C_p} = 0{,}810~\mathrm{mV/A} ,
  & V_{p,k+1} &= \frac{\Tc\,I}{C_p} = 8{,}10~\mathrm{mV} , \\[2pt]
  \alpha_1 &= -\kappa = -0{,}201 \cdot 10^{-3}~\mathrm{A}^{-1} ,
  & s_{k+1} &= 0{,}6 - \kappa\,I = 0{,}5980 .
\end{align*}

\paragraph{Lungo l'orizzonte.} La tabella riporta i quattro termini della
\eqref{eq:beta} ai passi 1, 2 e 30, con i valori da cui sono calcolati. Su
tutto l'orizzonte il SOC resta nel tratto tra $0{,}5$ e $0{,}75$, quindi
$R_p' = 0{,}0028~\Omega$ e $C_p' = 59636$~F non cambiano.

\begin{center}
\begin{tabular}{lrrr}
\toprule
 & $j = 1$ & $j = 2$ & $j = 30$ \\
\midrule
\multicolumn{4}{l}{\emph{Valori di inizio passo}} \\
$s_{k+j-1}$ & 0,6000 & 0,5980 & 0,5417 \\
$V_{p,k+j-1}$ [mV] & 0 & 8,10 & 17,27 \\
$\alpha_{j-1}$ [$10^{-3}~\mathrm{A}^{-1}$] & 0 & $-0{,}201$ & $-5{,}832$ \\
$\beta_{j-1}$ [mV/A] & 0 & 0,810 & 1,572 \\
\midrule
\multicolumn{4}{l}{\emph{Parametri letti in $s_{k+j-1}$}} \\
$C_p$ [F] & 24685 & 24565 & 21207 \\
$\tau$ [s] & 46,4 & 46,0 & 36,4 \\
$\tau'$ [s] & 181,2 & 180,6 & 161,8 \\
\midrule
\multicolumn{4}{l}{\emph{Termini della \eqref{eq:beta}} [$\mu$V/(A\,s)]} \\
$-\beta_{j-1}/\tau$ \quad ($V_p$ dipende da $I$) & 0 & $-17{,}6$ & $-43{,}2$ \\
$V_{p,k+j-1}\,\tau'\,\alpha_{j-1}/\tau^2$ \quad ($\tau$ tramite il SOC) & 0 & $-0{,}1$ & $-12{,}3$ \\
$1/C_p$ \quad ($I$ diretta) & 40,5 & 40,7 & 47,2 \\
$-I\,C_p'\,\alpha_{j-1}/C_p^2$ \quad ($C_p$ tramite il SOC) & 0 & $+0{,}2$ & $+7{,}7$ \\
Somma & 40,5 & 23,2 & $-0{,}6$ \\
\midrule
$\beta_j = \beta_{j-1} + \Tc \cdot \text{somma}$ [mV/A] & 0,810 & 1,274 & 1,560 \\
\bottomrule
\end{tabular}
\end{center}

\noindent
Tre cose si leggono nella tabella.
\begin{itemize}
  \item Al primo passo c'è solo il termine diretto, come già visto.
  \item Dal secondo passo compare il termine $-\beta_{j-1}/\tau$, di segno
        opposto, che frena la crescita di $\beta$: al passo 30 la somma è
        quasi nulla e $\beta$ si è assestata poco sotto il valore di $R_p$ a
        quel SOC ($1{,}72$~mV/A).
  \item I due termini che esistono solo perché i parametri dipendono dal SOC
        sono trascurabili all'inizio: al passo 2 valgono meno dell'1\% del
        termine diretto. Crescono però insieme ad $\alpha$, cioè con il numero
        di passi: al passo 30 valgono il 26\% e il 16\% del termine diretto,
        con segni opposti. La non linearità pesa sempre di più andando avanti
        lungo l'orizzonte.
\end{itemize}
Per la tensione ai morsetti i parametri si leggono nel SOC di fine passo, dove
$V_{oc}' = 0{,}8528$~V e $R_o' = 0{,}0004~\Omega$ su tutto il tratto. I quattro
termini della \eqref{eq:D} e i due coefficienti di sensibilità valgono:

\begin{center}
\begin{tabular}{lrrr}
\toprule
 & $j = 1$ & $j = 2$ & $j = 30$ \\
\midrule
\multicolumn{4}{l}{\emph{Termini della \eqref{eq:D}} [mV/A]} \\
$V_{oc}'\,\alpha_j$ \quad ($V_{oc}$ tramite il SOC) & $-0{,}172$ & $-0{,}343$ & $-5{,}145$ \\
$-\beta_j$ \quad ($V_p$ dipende da $I$) & $-0{,}810$ & $-1{,}274$ & $-1{,}560$ \\
$-R_o'\,\alpha_j\,I$ \quad ($R_o$ tramite il SOC) & $+0{,}001$ & $+0{,}002$ & $+0{,}024$ \\
$-R_o$ \quad ($I$ diretta) & $-8{,}239$ & $-8{,}238$ & $-8{,}216$ \\
\midrule
$d_j$ [mV/A] & $-9{,}220$ & $-9{,}853$ & $-14{,}897$ \\
$\alpha_j$ [$10^{-3}~\mathrm{A}^{-1}$] & $-0{,}201$ & $-0{,}402$ & $-6{,}033$ \\
\bottomrule
\end{tabular}
\end{center}

\noindent
In concreto: con 1~A in più di corrente la tensione predetta è più bassa di
$9{,}2$~mV dopo un passo e di $14{,}9$~mV dopo 30. Dopo un passo il coefficiente
è quasi solo la resistenza serie; lungo l'orizzonte cresce il contributo della
tensione a vuoto, perché la corrente in più ha scaricato di più la cella.

\subsubsection{Errore di discretizzazione}

Le derivate così calcolate sono esatte per il modello discreto. Rispetto a
quelle del modello continuo ereditano invece l'errore di Eulero, lo stesso
discusso nella sezione~\ref{sec:discreto} per lo stato: $\alpha_j$ coincide con
$\alpha(t_j)$, mentre $\beta_j$ e $d_j$ approssimano $\beta(t_j)$ e $d(t_j)$.
Raccogliendo $\beta_{j-1}$ nella \eqref{eq:beta},
\[
  \beta_j = \left(1 - \frac{\Tc}{\tau}\right)\beta_{j-1}
            + \Tc\left(\frac{V_{p,k+j-1}\,\tau'\,\alpha_{j-1}}{\tau^2}
            + \frac{1}{C_p} - \frac{I\,C_p'\,\alpha_{j-1}}{C_p^2}\right) ,
\]
compare lo stesso fattore $1 - \Tc/\tau$ che nella \eqref{eq:vp} moltiplica
$V_{p,k+j-1}$: la derivata ha la stessa dinamica dello stato, e quindi lo
stesso tipo di errore. Va ribadito che questo errore riguarda il modello di
predizione, non la linearizzazione: l'MPC ottimizza sul modello discreto, e le
derivate sono coerenti con quello.

\subsubsection{Dalle derivate al QP}

Con $d_j$ e $\alpha_j$ calcolati per ogni passo e per ogni cella, lo sviluppo
\eqref{eq:lin} è completo:
\[
  y_{k+j}(u_k) \simeq \bar Y_j + d_j\,\bigl(u_k - \bar u\bigr), \qquad
  s_{k+j}(u_k) \simeq \bar S_j + \alpha_j\,\bigl(u_k - \bar u\bigr), \qquad
  j = 1, \dots, p .
\]
Sono la \eqref{eq:c-y-lin} e la \eqref{eq:c-s-lin} negli istanti di
campionamento. Gli infiniti vincoli del tempo continuo diventano $p$ vincoli
per cella e l'integrale del costo una somma su $j$: è il QP della prossima
sezione. $\beta_j$ non vi compare: si esaurisce nel calcolo di $d_j$.

\section{Dalla linearizzazione al QP}\label{sec:qp}

La sezione~\ref{sec:derivate} fornisce, per ogni cella $n$ e per ogni passo
$j$, i quattro numeri della \eqref{eq:lin}: $\bar Y_{j,n}$, $\bar S_{j,n}$,
$d_{j,n}$ e $\alpha_{j,n}$. Qui si vede, in sintesi, come entrano nel QP; la
costruzione completa, con le motivazioni di ogni passaggio, è nel documento
\file{Problema_QP_MPC.pdf}.

\paragraph{L'ingresso lungo l'orizzonte.} Prima di costruire il QP conviene
precisare che cosa sia, matematicamente, l'incognita. All'istante $k$ l'MPC
predice l'evoluzione delle celle sui passi $k+1, \dots, k+p$, e in generale
potrebbe scegliere una corrente di bilanciamento diversa per ogni passo. Si
indica con
\[
  u^n_{k+j\,|\,k}, \qquad j = 0, 1, \dots, p ,
\]
la corrente che, all'istante $k$, si prevede di applicare alla cella $n$ al
passo $k+j$: il pedice prima della barra verticale dice \emph{quando agisce},
quello dopo \emph{quando è stata decisa}. È la notazione usuale dell'MPC. Allo
stesso modo $i_{k+j\,|\,k}$ è la corrente di pacco prevista per il passo
$k+j$. Con correnti che cambiano da un passo all'altro, il modello
\eqref{eq:soc}--\eqref{eq:y} andrebbe scritto con la corrente del passo,
$I^n_{k+j} = i_{k+j\,|\,k} + u^n_{k+j\,|\,k}$:
\begin{align*}
  s_{k+j} &= \min\bigl(s_{k+j-1} - \kappa^n I^n_{k+j-1},\; 1\bigr), \\
  V_{p,k+j} &= V_{p,k+j-1} + \Tc \left( -\frac{V_{p,k+j-1}}{\tau(s_{k+j-1})}
              + \frac{I^n_{k+j-1}}{C_p(s_{k+j-1})}\right), \\
  y_{k+j} &= V_{oc}(s_{k+j}) - V_{p,k+j} - R_o(s_{k+j})\, I^n_{k+j} .
\end{align*}
Lo stato all'istante $k+j$ è il risultato delle correnti dei passi
precedenti, da $k$ a $k+j-1$; la tensione ai morsetti dipende in più dalla
corrente del passo stesso, attraverso la caduta su $R_o$. In generale,
quindi, la tensione predetta è funzione di tutta la sequenza di correnti:
\begin{equation}
  y^n_{k+j} = y^n_{k+j}\bigl(u^n_{k\,|\,k},\; u^n_{k+1\,|\,k},\; \dots,\;
              u^n_{k+j\,|\,k}\bigr),
  \qquad j = 1, \dots, p .
  \tag{U1}\label{eq:u-sequenza}
\end{equation}
Il paper (eq.~6b--6c) non sceglie una sequenza: sceglie una sola corrente per
cella e la suppone costante su tutto l'orizzonte. Della corrente di pacco, di
cui non si ha un'anteprima, suppone lo stesso:
\begin{equation}
  u^n_{k+j\,|\,k} = u^n_{k\,|\,k} \;=:\; u^n_k ,
  \qquad
  i_{k+j\,|\,k} = i_k ,
  \qquad j = 1, \dots, p .
  \tag{U2}\label{eq:u-costante}
\end{equation}
Il simbolo $u^n_k$ usato in tutto il documento è proprio questa corrente
unica, decisa all'istante $k$. Con la \eqref{eq:u-costante} la corrente della
cella è la stessa a ogni passo, $I^n_{k+j} = i_k + u^n_k = I^n$, cioè la
costante della (1), e tutti gli argomenti della \eqref{eq:u-sequenza}
coincidono:
\begin{equation}
  y^n_{k+j} = y^n_{k+j}\bigl(u^n_k,\; u^n_k,\; \dots,\; u^n_k\bigr)
  \;\equiv\; y^n_{k+j}(u^n_k) ,
  \qquad j = 1, \dots, p .
  \tag{U3}\label{eq:u-unica}
\end{equation}
L'ultima scrittura è un leggero abuso di notazione: lo stesso simbolo
$y^n_{k+j}$ indica la funzione di un solo argomento che si ottiene ponendo
uguali tutti gli argomenti della \eqref{eq:u-sequenza}. È questa la funzione
che la \eqref{eq:lin} linearizza. Lo stesso vale per il SOC,
$s^n_{k+j}(u^n_k)$. Tre osservazioni sulla \eqref{eq:u-costante}:
\begin{itemize}
  \item \emph{È un'ipotesi della predizione, non del pacco.} Della corrente
        decisa in $k$ si applica solo il primo passo. All'istante $k+1$ il
        problema si risolve di nuovo, con le nuove misure, e si sceglie
        $u^n_{k+1} = u^n_{k+1\,|\,k+1}$, in generale diversa da
        $u^n_{k+1\,|\,k} = u^n_k$. La corrente applicata al pacco cambia
        quindi da un passo all'altro: è costante solo dentro ciascuna
        predizione.
  \item \emph{Riduce le incognite.} Invece di una sequenza di correnti per
        cella, le incognite sono $N$, una per cella. Il paper la giustifica
        osservando che le correnti di bilanciamento variano lentamente
        (sua Fig.~2).
  \item \emph{Rende la predizione una funzione di una variabile.} Poiché ogni
        $y^n_{k+j}$ dipende da un solo numero, la \eqref{eq:lin} è una retta
        in $u^n_k$, e non uno sviluppo in $j+1$ variabili.
\end{itemize}
Da qui in avanti si scrive solo $u^n_k$, e anche nelle predizioni si omette
``$|\,k$'': $y^n_{k+j}$ e $s^n_{k+j}$ sono sempre le grandezze previste
all'istante $k$.

\paragraph{Risposta libera e risposta forzata.} Conviene anzitutto isolare
nella \eqref{eq:lin} la parte che non dipende da $u^n_k$:
\begin{equation}
  y^n_{k+j}(u^n_k) \simeq Y^\lib_{j,n} + d_{j,n}\,u^n_k,
  \qquad Y^\lib_{j,n} = \bar Y_{j,n} - d_{j,n}\,\bar u^n .
  \label{eq:y0}
\end{equation}
$Y^\lib_{j,n}$ è la \emph{risposta libera}: il valore che la retta della cella
$n$ dà in $u^n_k = 0$. Il prodotto $d_{j,n}\,u^n_k$ è la \emph{risposta forzata},
l'effetto della corrente di bilanciamento.

\paragraph{Il problema quadratico.} Le incognite del QP sono raccolte in
\[
  z = \begin{bmatrix} u_k \\ \varepsilon \\ \varepsilon_y \end{bmatrix},
\]
con $u_k \in \mathbb{R}^N$ le correnti di bilanciamento scelte all'istante $k$, $\varepsilon$ le
variabili ausiliarie del costo (nessuna per $J_t$, $p$ per $J_m$, $2p$ per
$J_\Delta$) e $\varepsilon_y \ge 0$ la variabile che rilassa il vincolo di
tensione (eq.~17). Il problema è
\[
  \min_z \; \tfrac12\, z^\top H z + f^\top z
  \qquad \text{con} \qquad A z \le b, \quad A_{eq}\, z = b_{eq} .
\]
Qui $f$ è il vettore del termine lineare del costo, secondo la notazione di
\texttt{quadprog}: non è la funzione $f(\zeta, I)$ della
sezione~\ref{sec:modello}. In tutte le formulazioni il costo contiene
$\rho\, u_k^\top u_k + W \varepsilon_y^2$, dove $\rho$ è il peso sulle correnti di
bilanciamento (la matrice $R$ del paper vale $\rho$ volte l'identità) e $W$
il peso su $\varepsilon_y$. Il peso $\rho$ è scalato con l'orizzonte,
$\rho = \rho_{\mathrm{rif}}\, p / p_{\mathrm{rif}}$, dove $\rho_{\mathrm{rif}}$
è \texttt{par.mpc.Jt}, \texttt{par.mpc.Jm} o \texttt{par.mpc.Jdelta} secondo
la formulazione. I due termini contribuiscono a $H$ con $2\rho$ sulla
diagonale delle correnti e $2W$ su $\varepsilon_y$; da dove venga il fattore 2
si vede nella sezione~\ref{sec:costo}, a proposito di $J_t$.

\subsection{Costo}\label{sec:costo}

Il costo è la somma di un termine di bilanciamento, diverso per ogni
formulazione, e dei due termini comuni $\rho\, u_k^\top u_k$ e $W \varepsilon_y^2$. Il termine
di $J_t$ è svolto qui passo per passo: è l'unico in cui la linearizzazione
produce termini di secondo grado, e i passaggi sono gli stessi che servono
per leggere $H$ e $f$ da qualunque costo. Quelli di $J_m$ e $J_\Delta$ sono
di primo grado e si leggono subito.

\subsubsection{$J_t$, inseguimento della cella nominale (eq.~7)}

\paragraph{Il costo da minimizzare.} La eq.~7 del paper somma, passo per
passo, il quadrato della lunghezza di un vettore di $N$ componenti, una per
cella. Scritta elemento per elemento, con la tensione come grandezza
bilanciata, è
\[
  J_t(u_k) \;=\; \sum_{j=1}^{p} \sum_{n=1}^{N} \bigl(y^n_{k+j} - y^0_{k+j}\bigr)^2
  \;+\; \rho \sum_{n=1}^{N} (u^n_k)^2 ,
\]
con:
\begin{itemize}
  \item $y^n_{k+j}$, la tensione predetta della cella $n$ all'istante $k+j$:
        per la \eqref{eq:u-unica} dipende dalla sola incognita $u^n_k$;
  \item $y^0_{k+j}$, la tensione della \emph{cella nominale} all'istante
        $k+j$, cioè il riferimento che ogni cella insegue. Come nel paper,
        l'apice 0 al posto dell'indice di cella $n$ indica la cella nominale:
        parametri senza deviazioni, partita dalla media degli stati del
        pacco. La calcola il blocco \file{Cella_nominale}. È la stessa per
        tutte le celle e non dipende da $u_k$: per il QP è un numero noto. Il
        paper la scrive $\sigma^0$, che con $\sigma = y$ è proprio $y^0$; nella
        eq.~7 compare come $\sigma^0_{k+1}$, ma la tensione nominale cambia
        lungo l'orizzonte e va letta $\sigma^0_{k+j}$;
  \item $\rho$, il peso sulle correnti.
\end{itemize}
Il primo termine misura quanto ogni cella si allontana dalla cella nominale,
a ogni passo; il secondo rende costose le correnti grandi.

\paragraph{Due tensioni da non confondere.} $y^0_{k+j}$ e $Y^\lib_{j,n}$ sono
due tensioni diverse. $y^0_{k+j}$ è della cella nominale, una sola per tutto
il pacco. $Y^\lib_{j,n}$, definita nella \eqref{eq:y0}, è della cella $n$: è
il valore che la sua retta dà in $u^n_k = 0$. Coincide con la tensione che la
cella avrebbe senza bilanciamento solo se $\bar u^n = 0$; altrimenti è il
prolungamento della tangente fino a $u^n_k = 0$.

\paragraph{Passo 1: sostituire la retta.} Con la \eqref{eq:y0} la distanza
dalla cella nominale diventa
\[
  y^n_{k+j} - y^0_{k+j}
  = Y^\lib_{j,n} + d_{j,n}\,u^n_k - y^0_{k+j}
  = d_{j,n}\,u^n_k + c_{j,n} ,
  \qquad
  c_{j,n} = Y^\lib_{j,n} - y^0_{k+j} .
\]
$c_{j,n}$ è lo \emph{scarto}: di quanto, secondo la retta, la cella $n$
starebbe sopra ($c_{j,n} > 0$) o sotto ($c_{j,n} < 0$) la cella nominale
all'istante $k+j$ con $u^n_k = 0$. Come $d_{j,n}$, è un numero noto prima di
ottimizzare. Il costo diventa
\begin{equation}
  J_t(u_k) = \sum_{j=1}^{p} \sum_{n=1}^{N} \bigl(d_{j,n}\,u^n_k + c_{j,n}\bigr)^2
           + \rho \sum_{n=1}^{N} (u^n_k)^2 .
  \tag{Q1}\label{eq:q-jt}
\end{equation}

\paragraph{Passo 2: separare le celle.} Le due somme sono finite, quindi si
può scambiarne l'ordine e portare il termine in $\rho$ dentro la somma sulle
celle:
\begin{equation}
  J_t(u_k) = \sum_{n=1}^{N} J^n(u^n_k) ,
  \qquad
  J^n(u^n_k) = \sum_{j=1}^{p} \bigl(d_{j,n}\,u^n_k + c_{j,n}\bigr)^2 + \rho\,(u^n_k)^2 .
  \tag{Q2}\label{eq:q-separa}
\end{equation}
Ogni $J^n$ contiene una sola incognita, $u^n_k$: il costo è la somma di $N$
funzioni di una variabile, una per cella. Ne segue che nel costo non
compaiono prodotti misti $u^n_k u^m_k$ con $n \ne m$, e quindi che il blocco di
$H$ relativo alle correnti è diagonale. Le celle restano accoppiate solo dai
vincoli, attraverso $\sum_n u^n_k = 0$ (sezione~\ref{sec:vincoli}).

\paragraph{Passo 3: sviluppare il quadrato.} Si fissa una cella. Per
alleggerire la scrittura, in questo passo e nel prossimo l'indice $n$ è
sottinteso: $d_j$, $c_j$ e $u$ stanno per $d_{j,n}$, $c_{j,n}$ e $u^n_k$. Ogni
addendo della \eqref{eq:q-separa} è il quadrato di un binomio:
\[
  (d_j\,u + c_j)^2 = d_j^2\,u^2 + 2\,d_j\,c_j\,u + c_j^2 .
\]
Con un orizzonte di due passi, $p = 2$, il conto per esteso è
\begin{align*}
  J^n(u) &= (d_1 u + c_1)^2 + (d_2 u + c_2)^2 + \rho\,u^2 \\
         &= d_1^2 u^2 + 2 d_1 c_1 u + c_1^2
            \;+\; d_2^2 u^2 + 2 d_2 c_2 u + c_2^2
            \;+\; \rho\,u^2 \\
         &= \bigl(d_1^2 + d_2^2 + \rho\bigr)\, u^2
            + 2\,\bigl(d_1 c_1 + d_2 c_2\bigr)\, u
            + \bigl(c_1^2 + c_2^2\bigr) .
\end{align*}
Nell'ultima riga si sono raccolti i termini con $u^2$, quelli con $u$ e
quelli senza $u$. Per un orizzonte qualsiasi si raccoglie allo stesso modo:
\begin{equation}
  J^n(u) = \underbrace{\Bigl(\sum_{j=1}^{p} d_j^2 + \rho\Bigr)}_{\text{coefficiente di } u^2}\, u^2
         + \underbrace{\Bigl(2\sum_{j=1}^{p} d_j\,c_j\Bigr)}_{\text{coefficiente di } u}\, u
         + \underbrace{\sum_{j=1}^{p} c_j^2}_{\text{costante}} .
  \tag{Q3}\label{eq:q-sviluppo}
\end{equation}
$J^n$ è quindi un polinomio di secondo grado in $u$, con tre coefficienti
noti.

\paragraph{Passo 4: confrontare con la forma di \texttt{quadprog}.}
\texttt{quadprog} vuole il costo nella forma $\tfrac12\, z^\top H z + f^\top
z$. Le correnti occupano le prime $N$ posizioni di $z$, quindi $u^n_k = z_n$,
e i suoi coefficienti in $H$ e in $f$ sono $H_{nn}$ e $f_n$. Poiché $H$ non ha
prodotti misti tra le correnti (passo 2), la parte di questa espressione che
contiene $u^n_k$ è
\[
  \tfrac12\, H_{nn}\,(u^n_k)^2 + f_n\, u^n_k
\]
(è la regola di \file{Problema_QP_MPC.pdf}, sezione~2.2). Perché i due costi
abbiano lo stesso punto di minimo devono avere gli stessi coefficienti di
$u^2$ e di $u$; la costante si può trascurare, perché sommare una costante a
un costo ne sposta il valore ma non il punto in cui è minimo. Confrontando con
la \eqref{eq:q-sviluppo}:
\[
  \tfrac12\, H_{nn} = \sum_{j=1}^{p} d_j^2 + \rho
  \quad\Longrightarrow\quad
  H_{nn} = 2\sum_{j=1}^{p} d_j^2 + 2\rho ,
  \qquad\qquad
  f_n = 2\sum_{j=1}^{p} d_j\,c_j .
\]
Rimettendo l'indice $n$ e ricordando che $c_{j,n} = Y^\lib_{j,n} - y^0_{k+j}$:
\begin{equation}
  H_{nn} = 2\sum_{j=1}^{p} d_{j,n}^2 + 2\rho, \qquad
  f_n = 2\sum_{j=1}^{p} d_{j,n}\,\bigl(Y^\lib_{j,n} - y^0_{k+j}\bigr) .
  \label{eq:jt-hf}
\end{equation}
I due fattori 2 hanno origini diverse. Quello in $H$ compensa il
$\tfrac12$ della forma di \texttt{quadprog}, che è una convenzione: con esso il
gradiente del costo è $Hz + f$ e $H$ è la matrice delle derivate seconde.
Quello in $f$ viene dal doppio prodotto del binomio. Lo stesso ragionamento,
applicato a $\rho\,(u^n_k)^2$ e a $W \varepsilon_y^2$ da soli, dà i contributi $2\rho$ e
$2W$ citati all'inizio della sezione.

\paragraph{Verifica: il minimo per una cella.} Ignorando per un momento i
vincoli, $J^n$ è una parabola in $u$. Il coefficiente di $u^2$,
$\sum_j d_j^2 + \rho$, è positivo, quindi la parabola è rivolta verso l'alto e
ha un solo minimo, dove la derivata si annulla:
\[
  \frac{dJ^n}{du} = 2\Bigl(\sum_j d_j^2 + \rho\Bigr) u + 2\sum_j d_j\,c_j
                  = H_{nn}\, u + f_n = 0 .
\]
Ne segue
\begin{equation}
  u^{n,*}_k = -\frac{f_n}{H_{nn}}
        = \frac{\sum_{j} d_{j,n}\,\bigl(y^0_{k+j} - Y^\lib_{j,n}\bigr)}
               {\sum_{j} d_{j,n}^2 + \rho} .
  \tag{Q4}\label{eq:q-minimo}
\end{equation}
Lo stesso risultato si ottiene derivando direttamente la
\eqref{eq:q-separa}, senza sviluppare i quadrati: per la regola della
funzione composta $\frac{d}{du}(d_j u + c_j)^2 = 2 d_j\,(d_j u + c_j)$, e
sommando su $j$ si ritrova la stessa equazione. È una conferma che
$H_{nn}$ e $f_n$ sono giusti.

La \eqref{eq:q-minimo} è la soluzione di un problema ai minimi quadrati: ci
sono $p$ scarti da colmare e una sola corrente, che al passo $j$ sposta la
tensione di $d_{j,n}\,u^n_k$. Non potendo annullarli tutti con un solo numero,
si sceglie quello che rende minima la somma dei quadrati di ciò che resta. Il
peso $\rho$ al denominatore frena la corrente. Anche il segno torna: se la
cella sta sopra la cella nominale ($c_{j,n} > 0$), essendo $d_{j,n} < 0$
risulta $u^{n,*}_k > 0$, cioè la cella viene scaricata di più e la sua tensione
scende verso $y^0_{k+j}$.

\paragraph{Esempio numerico.} Si considera una cella sola e un orizzonte di
due passi, $p = 2$; come nei passi 3 e 4, l'indice di cella è sottinteso, $u$
sta per $u^n_k$ e $u^*$ per $u^{n,*}_k$.
I coefficienti di sensibilità sono quelli calcolati
nell'esempio della sezione~\ref{sec:derdiscreto} (cella con fattori di
deviazione uguali a 1, $s_k = 0{,}6$, $I = 10$~A), nelle colonne $j = 1$ e
$j = 2$:
\[
  d_1 = -9{,}220 \cdot 10^{-3}~\text{V/A}, \qquad
  d_2 = -9{,}853 \cdot 10^{-3}~\text{V/A}.
\]
Gli scarti $c_j = Y^\lib_j - y^0_{k+j}$ sono valori di prova,
$c_1 = +20$~mV e $c_2 = +25$~mV: la cella sta sopra la cella nominale e se ne
allontana. Il peso, con i valori di
\file{init_parametri.m} (\texttt{par.mpc.Jt} $= 2 \cdot 10^{-4}$,
$p_{\mathrm{rif}} = 5$), è $\rho = 2 \cdot 10^{-4} \cdot 2/5 = 8 \cdot
10^{-5}$. Il costo è in V$^2$ e $u$ in ampere, quindi $d_j^2$ e $\rho$ sono in
V$^2$/A$^2$ e $d_j c_j$ in V$^2$/A.

I due coefficienti della \eqref{eq:q-sviluppo} che contengono $u$ valgono
\begin{align*}
  d_1^2 + d_2^2 + \rho
    &= 0{,}8501 \cdot 10^{-4} + 0{,}9708 \cdot 10^{-4} + 0{,}8000 \cdot 10^{-4}
     = 2{,}6209 \cdot 10^{-4} , \\
  2\,(d_1 c_1 + d_2 c_2)
    &= 2\,\bigl(-1{,}8440 \cdot 10^{-4} - 2{,}4633 \cdot 10^{-4}\bigr)
     = -8{,}6145 \cdot 10^{-4} ,
\end{align*}
quindi $H_{11} = 2 \cdot 2{,}6209 \cdot 10^{-4} = 5{,}2418 \cdot 10^{-4}$ e
$f_1 = -8{,}6145 \cdot 10^{-4}$. La \eqref{eq:q-minimo} dà
\[
  u^* = -\frac{f_1}{H_{11}} = \frac{8{,}6145}{5{,}2418} = 1{,}643~\text{A}.
\]
Con questa corrente la retta dà, rispetto alla cella nominale,
\begin{align*}
  y_{k+1} - y^0_{k+1} &= d_1 u^* + c_1 = -15{,}15 + 20 = +4{,}85~\text{mV}
    && \text{al primo passo}, \\
  y_{k+2} - y^0_{k+2} &= d_2 u^* + c_2 = -16{,}19 + 25 = +8{,}81~\text{mV}
    && \text{al secondo}:
\end{align*}
la cella si avvicina alla cella nominale ma non la raggiunge, perché $\rho$
rende la corrente costosa.
Con $\rho = 0$ si avrebbe $u^* = 4{,}3073/1{,}8209 = 2{,}365$~A e gli scarti
residui diventerebbero $-1{,}81$~mV e $+1{,}69$~mV: uno sotto e uno sopra,
perché una sola corrente non può annullare due scarti diversi e sceglie il
compromesso con la minima somma dei quadrati.

\paragraph{Più celle e vincoli.} Con $N$ celle ogni cella porta i propri
$H_{nn}$ e $f_n$ nella propria posizione, e il blocco delle correnti in $H$ è
la matrice diagonale con elementi $H_{11}, \dots, H_{NN}$. Le correnti ideali
\eqref{eq:q-minimo} in generale non rispettano $\sum_n u^n_k = 0$: se per
esempio tutte le celle stanno sopra la cella nominale, sono tutte positive. Il
vincolo sulla somma lega le celle e \texttt{quadprog} cerca, tra le correnti a
somma nulla e dentro i limiti, quella di costo minimo. La parte dello scarto
comune a tutte le celle non si può correggere: il bilanciatore sposta carica
da una cella all'altra, non ne toglie al pacco.

\subsubsection{$J_m$, massimizzazione della cella più bassa (eq.~9--10)}

Il termine di bilanciamento è $-\sum_j \varepsilon_j$. Le $\varepsilon_j$
compaiono solo al primo grado, con coefficiente $-1$: per la regola del
passo~4 danno $f_{N+j} = -1$ e nessun contributo a $H$. Il minimo sulle celle
è espresso dai vincoli $\varepsilon_j \le y^n_{k+j}$ per ogni $n$: minimizzare
$-\sum_j \varepsilon_j$ spinge ogni $\varepsilon_j$ verso l'alto finché tocca
la tensione della cella più bassa al passo $j$.

\subsubsection{$J_\Delta$, minima differenza tra le celle (eq.~12--13)}

Il termine di bilanciamento è $\sum_j \varepsilon_{p+j} - \sum_j
\varepsilon_j$, ancora di primo grado: $f_{N+j} = -1$ e $f_{N+p+j} = +1$,
nessun contributo a $H$. Ai vincoli di $J_m$ si aggiungono
$\varepsilon_{p+j} \ge y^n_{k+j}$, che spingono ogni $\varepsilon_{p+j}$ verso
il basso fino alla cella più alta.

\subsection{Vincoli}\label{sec:vincoli}

Sostituendo la predizione linearizzata, tutti i vincoli diventano lineari in
$z$ (per ogni $j = 1,\dots,p$ e $n = 1,\dots,N$):
\begin{align}
  -d_{j,n}\,u^n_k - \varepsilon_y &\le Y^\lib_{j,n} - y_{\min}
     && \text{tensione minima rilassata, eq.~(6e) e (17)} \\
  \varepsilon_j - d_{j,n}\,u^n_k &\le Y^\lib_{j,n}
     && \text{eq.~(10) e (13a), per } J_m \text{ e } J_\Delta \\
  d_{j,n}\,u^n_k - \varepsilon_{p+j} &\le -Y^\lib_{j,n}
     && \text{eq.~(13b), per } J_\Delta \\
  u_{\min} \le u^n_k &\le u_{\max}
     && \text{eq.~(6d)} \\
  -\varepsilon_y &\le 0 \\
  \textstyle\sum_{n} u^n_k &= 0
     && \text{eq.~(6f)}
\end{align}
Se la grandezza bilanciata è il SOC anziché la tensione, nel costo e nei
vincoli sulle $\varepsilon$ si usano $\bar S$ e $\alpha$ al posto di $\bar Y$ e
$d$; il vincolo di tensione minima usa sempre $\bar Y$ e $d$.

\subsection{Punto di partenza del solutore}

Nel punto $u_k = \bar u$ la predizione linearizzata coincide con la traiettoria
di linearizzazione, quindi un punto ammissibile si scrive direttamente:
\[
  u_k = \bar u, \qquad
  \varepsilon_j = \min_n \bar Y_{j,n}, \qquad
  \varepsilon_{p+j} = \max_n \bar Y_{j,n}, \qquad
  \varepsilon_y = \max\bigl(0,\; y_{\min} - \min_{j,n} \bar Y_{j,n}\bigr).
\]
\texttt{quadprog} (algoritmo \texttt{active-set}) parte da lì. Se non trova una
soluzione, l'uscita è questo punto: la mossa precedente viene mantenuta.

\section{Riepilogo di un passo di controllo}\label{sec:riepilogo}

\begin{enumerate}
  \item Si campionano $s_k$, $V_{p,k}$ di ogni cella e la corrente di pacco $i_k$.
  \item \file{Cella_nominale}: si calcolano le tensioni della cella nominale,
        $y^0_{k+1},\dots,y^0_{k+p}$, integrando una cella senza deviazioni
        dalla media degli stati.
  \item \file{Predizione_linearizzata}: con la corrente $i_k + u^n_{k-1}$ si
        calcolano la traiettoria di linearizzazione $\bar Y$, $\bar S$, con le
        \eqref{eq:soc}--\eqref{eq:y}, e i coefficienti di sensibilità $d$ e
        $\alpha$, con le \eqref{eq:beta}--\eqref{eq:D}.
  \item \file{Problema_QP}: si costruiscono $H$, $f$, $A$, $b$, $A_{eq}$,
        $b_{eq}$ e il punto di partenza.
  \item \file{Solutore_QP}: \texttt{quadprog} restituisce $z$; al pacco si
        applicano solo le prime $N$ componenti, $u_k$.
  \item $u_k$ viene memorizzata: sarà il punto di linearizzazione del passo
        successivo.
\end{enumerate}
La linearizzazione viene fatta una sola volta per passo, senza iterare.

\section{Quanto è accurata}\label{sec:accuratezza}

Nel documento compaiono due approssimazioni diverse, che conviene tenere
distinte. La prima è la discretizzazione del modello
(sezione~\ref{sec:discreto}): riguarda il modello di predizione, ed è la
stessa per il problema non lineare e per il QP. La seconda è la
linearizzazione, cioè la sostituzione della predizione con la retta
\eqref{eq:lin}: è la sola differenza tra il QP e il problema non lineare, ed è
quella di cui si occupa questa sezione.

La linearizzazione introduce errore solo dove il modello non è lineare nella
corrente:
\begin{itemize}
  \item il SOC è esattamente lineare in $u_k$, come mostra la \eqref{eq:alpha},
        tranne quando una cella satura a 1;
  \item la tensione non lo è per i prodotti $R_o(s)\,I$ e $I/C_p(s)$, per
        $\tau(s)$ e per i cambi di pendenza delle tabelle quando il SOC
        attraversa uno dei 7 punti lungo l'orizzonte.
\end{itemize}
L'errore cresce con la distanza tra $u_k$ e $u_{k-1}$ e con la lunghezza
dell'orizzonte $p\,\Tc$: come si è visto nell'esempio numerico, i termini
dovuti alla non linearità crescono con il numero di passi.

La Tabella~\ref{tab:confronto} confronta, sullo scenario a f.e.m.\ costante, il
QP linearizzato con la soluzione del problema non lineare (\texttt{nlmpc}),
usando un unico passo per integrazione e controllo. In una prova separata, con
$\Tc = 20$~s e $p = 30$, ripetere la linearizzazione tre volte per passo non ha
cambiato le durate.

\begin{table}[ht]
\centering
\begin{tabular}{llcccc}
\toprule
 & & \multicolumn{2}{c}{Durata [s]} & Scarto massimo & Sforzo medio $\bar w$ [A$^2$] \\
Configurazione & Costo & \texttt{nlmpc} & QP & su $u_k$ [A] & \texttt{nlmpc} / QP \\
\midrule
$N=5$, $\Tc=1$ s, $p=5$   & $J_t$      & 2529 & 2529 & $3{,}7\cdot10^{-3}$ & 5,3526 / 5,3526 \\
(orizzonte 5 s)           & $J_m$      & 2478 & 2478 & $9{,}3\cdot10^{-4}$ & 0,7825 / 0,7825 \\
                          & $J_\Delta$ & 2488 & 2488 & $3{,}5\cdot10^{-3}$ & 1,2847 / 1,2847 \\
\midrule
$N=3$, $\Tc=20$ s, $p=30$ & $J_t$      & 8120 & 8120 & $0{,}17$            & 1,0733 / 1,0742 \\
(orizzonte 600 s)         & $J_m$      & 8120 & 8120 & $8{,}2\cdot10^{-3}$ & 1,1376 / 1,1375 \\
                          & $J_\Delta$ & 8120 & 8120 & $0{,}28$            & 1,2143 / 1,2128 \\
\bottomrule
\end{tabular}
\caption{QP linearizzato a confronto con il problema non lineare. Gli scarti
maggiori del secondo caso compaiono in pochi passi isolati. Lo sforzo di
bilanciamento al passo $k$ è $w_k = u_k^\top u_k = \sum_n (u^n_k)^2$ (nel
paper $e_k$), e $\bar w = \frac1K \sum_k w_k$ è la sua media sui $K$ passi
della simulazione.}
\label{tab:confronto}
\end{table}

\section{Dove si trova nel codice}\label{sec:codice}

Nel modello \file{Bilanciamento_MPC.slx} ogni passaggio è un blocco del
sottosistema \file{Controllo}, e ogni blocco chiama una funzione:

\begin{center}
\small
\begin{tabular}{lll}
\toprule
Passaggio & Blocco in \file{Controllo} & File \\
\midrule
Cella nominale $y^0_{k+j}$ & \file{Cella_nominale} & \file{cella_nominale_semplice.m} \\
Predizione e derivate, sezione~\ref{sec:derdiscreto} & \file{Predizione_linearizzata} & \file{qp_predizione.m} \\
Valore e derivata dei parametri, \eqref{eq:tabella} & --- & \file{qp_tabella.m} \\
Matrici del QP e punto di partenza & \file{Problema_QP} & \file{qp_problema.m} \\
Soluzione & \file{Solutore_QP} & \file{qp_solutore.m} \\
\bottomrule
\end{tabular}
\end{center}

\noindent
La costruzione delle matrici del QP ha un documento a parte,
\file{Problema_QP_MPC.pdf}, e una versione da leggere,
\file{problema_qp_semplice.m}. \file{qp_predizione.m} applica le formule della
predizione a tutte le celle insieme, con vettori. Per leggerle con più agio ci
sono tre file che il modello non usa:
\begin{itemize}
  \item \file{predizione_linearizzata_semplice.m} fa lo stesso calcolo una
        cella alla volta, con numeri singoli, ed è diviso nelle stesse parti di
        questo documento (\file{qp_cella_nominale.m} è la versione con vettori
        del calcolo della cella nominale);
  \item \file{Predizione_linearizzata_blocchi.slx} fa lo stesso calcolo con
        blocchi Simulink, un sottosistema per ogni parte;
  \item \file{esempio_sensibilita_ode45.m} applica le formule generali della
        sezione~\ref{sec:sens} a un sistema massa-molla-smorzatore.
\end{itemize}
\file{predizione_linearizzata_semplice.m} usa la notazione di questo documento:
ogni variabile ha il nome del simbolo corrispondente, e la derivata rispetto a
$s$, qui indicata con l'apice, nel codice ha il prefisso \file{ds_}, che
nomina anche la variabile di derivazione.

\begin{center}
\begin{tabular}{ll@{\qquad}ll}
\toprule
Simbolo & Nome nel codice & Simbolo & Nome nel codice \\
\midrule
$s_k$, $V_{p,k}$ & \file{s_k}, \file{Vp_k} & $s_{k+j}$, $V_{p,k+j}$ & \texttt{s}, \texttt{Vp} \\
$i_k$, $\bar u$ & \file{i_pacco_k}, \file{u_bar} & $I$ & \file{I_cella} \\
$\Tc$, $\kappa$ & \texttt{Tc}, \texttt{kappa} & $\alpha_j$, $\beta_j$ & \file{alfa_j}, \file{beta_j} \\
$V_{oc}$, $R_o$ & \texttt{Voc}, \texttt{Ro} & $V_{oc}'$, $R_o'$ & \file{ds_Voc}, \file{ds_Ro} \\
$R_p$, $C_p$, $\tau$ & \texttt{Rp}, \texttt{Cp}, \texttt{tau} & $R_p'$, $C_p'$, $\tau'$ & \file{ds_Rp}, \file{ds_Cp}, \file{ds_tau} \\
$\varphi_{R_o}$, $R_o^{\nom}$ & \file{phi_Ro}, \file{Ro_tab} & $y_{k+j}$, $d_j$ & \texttt{y}, \file{d_j} \\
\bottomrule
\end{tabular}
\end{center}

\noindent
I fattori e le tabelle degli altri parametri seguono lo stesso schema
(\file{phi_Voc}, \file{Voc_tab}, e così via). Le quattro matrici restituite
si chiamano \file{Y_bar}, \file{S_bar}, \texttt{d}, \texttt{alfa}: raccolgono
$y$, $s$, $d$ e $\alpha$ di tutti i passi e di tutte le celle.

\end{document}
